%%%%%%%%%%%%%%%%%%%%%%%%%%%%%%%%%%%%%%%%%%%%%%%%%%%%%%%%%%%%
% 8.7 REGRESSION
% The Composition of Advanced Mathematics
%%%%%%%%%%%%%%%%%%%%%%%%%%%%%%%%%%%%%%%%%%%%%%%%%%%%%%%%%%%%

\section{Regression}
\label{sec:regression}

\prerequisite{
Descriptive statistics, covariance, correlation, sampling distributions,
statistical estimation, confidence intervals, hypothesis testing,
linear algebra, matrix multiplication, derivatives, optimization, and
basic probability.
}

\learningobjectives{
By the end of this section, the reader should be able to:
\begin{itemize}
    \item explain the purpose of regression modeling;
    \item distinguish response and explanatory variables;
    \item formulate a simple linear regression model;
    \item distinguish population regression parameters from estimated
          regression coefficients;
    \item derive ordinary least squares estimators;
    \item interpret slope and intercept coefficients;
    \item compute fitted values and residuals;
    \item decompose variability using sums of squares;
    \item define \(R^2\);
    \item connect correlation and simple linear regression;
    \item understand assumptions of the classical linear model;
    \item derive sampling distributions of OLS estimators under normal
          errors;
    \item construct confidence intervals for regression coefficients;
    \item perform \(t\)-tests for regression coefficients;
    \item construct confidence intervals for the mean response;
    \item construct prediction intervals for future observations;
    \item distinguish confidence and prediction intervals;
    \item formulate multiple linear regression;
    \item express regression in matrix form;
    \item derive the normal equations;
    \item derive the OLS estimator
          \((X^TX)^{-1}X^Ty\);
    \item interpret partial regression coefficients;
    \item understand interaction terms;
    \item understand polynomial regression;
    \item work with indicator variables;
    \item understand categorical predictors;
    \item understand adjusted \(R^2\);
    \item interpret residual standard error;
    \item diagnose nonlinearity, heteroskedasticity, and outliers;
    \item understand leverage and influence;
    \item define Cook's distance conceptually;
    \item understand multicollinearity;
    \item interpret variance inflation factors;
    \item distinguish interpolation from extrapolation;
    \item understand omitted-variable bias conceptually;
    \item distinguish association from causation in regression;
    \item understand robust standard errors conceptually;
    \item recognize weighted least squares;
    \item recognize generalized least squares;
    \item understand model selection conceptually;
    \item recognize ridge and lasso regression;
    \item understand logistic regression as a generalized linear model
          preview;
    \item connect regression with ANOVA, experimental design,
          statistical learning, and mathematical modeling.
\end{itemize}
}

%%%%%%%%%%%%%%%%%%%%%%%%%%%%%%%%%%%%%%%%%%%%%%%%%%%%%%%%%%%%
\subsection{What Is Regression?}
\label{subsec:what-is-regression}
%%%%%%%%%%%%%%%%%%%%%%%%%%%%%%%%%%%%%%%%%%%%%%%%%%%%%%%%%%%%

Regression studies relationships between a response variable and one or
more explanatory variables.

Suppose the response is denoted

\[
Y
\]

and explanatory variables are denoted

\[
X_1,\ldots,X_p.
\]

The central goal is to model how the conditional distribution of \(Y\)
changes as the explanatory variables change.

A regression model may be used for:

\[
\boxed{
\text{description},
}
\]

\[
\boxed{
\text{prediction},
}
\]

\[
\boxed{
\text{estimation},
}
\]

and, under stronger assumptions,

\[
\boxed{
\text{causal analysis}.
}
\]

%%%%%%%%%%%%%%%%%%%%%%%%%%%%%%%%%%%%%%%%%%%%%%%%%%%%%%%%%%%%
\subsection{Response and Explanatory Variables}
\label{subsec:response-explanatory-regression}
%%%%%%%%%%%%%%%%%%%%%%%%%%%%%%%%%%%%%%%%%%%%%%%%%%%%%%%%%%%%

\begin{definitionbox}{Response Variable}
The response variable is the outcome whose behavior is being modeled or
predicted.
\end{definitionbox}

\begin{definitionbox}{Explanatory Variable}
An explanatory variable is a variable used to explain, predict, or
describe changes in the response.
\end{definitionbox}

Alternative names include:

\[
\boxed{
\text{predictor},
}
\]

\[
\boxed{
\text{covariate},
}
\]

or

\[
\boxed{
\text{independent variable}.
}
\]

The term ``independent variable'' does not necessarily imply
probabilistic independence.

%%%%%%%%%%%%%%%%%%%%%%%%%%%%%%%%%%%%%%%%%%%%%%%%%%%%%%%%%%%%
\subsection{Simple Linear Regression}
\label{subsec:simple-linear-regression}
%%%%%%%%%%%%%%%%%%%%%%%%%%%%%%%%%%%%%%%%%%%%%%%%%%%%%%%%%%%%

The simplest regression model uses one explanatory variable.

\begin{definitionbox}{Simple Linear Regression Model}
The simple linear regression model is

\[
\boxed{
Y_i
=
\beta_0
+
\beta_1X_i
+
\varepsilon_i,
}
\]

for

\[
i=1,\ldots,n.
\]
\end{definitionbox}

Here:

\[
\boxed{
\beta_0
=
\text{intercept},
}
\]

\[
\boxed{
\beta_1
=
\text{slope},
}
\]

and

\[
\boxed{
\varepsilon_i
=
\text{random error}.
}
\]

%%%%%%%%%%%%%%%%%%%%%%%%%%%%%%%%%%%%%%%%%%%%%%%%%%%%%%%%%%%%
\subsection{Conditional Mean Interpretation}
\label{subsec:conditional-mean-regression}
%%%%%%%%%%%%%%%%%%%%%%%%%%%%%%%%%%%%%%%%%%%%%%%%%%%%%%%%%%%%

If

\[
\mathbb{E}[\varepsilon_i\mid X_i]=0,
\]

then

\[
\boxed{
\mathbb{E}[Y_i\mid X_i]
=
\beta_0+\beta_1X_i.
}
\]

Thus the regression line describes the conditional mean response.

The model does not state that every observation lies exactly on the line.

Instead,

\[
\boxed{
Y_i
=
\text{systematic component}
+
\text{random deviation}.
}
\]

%%%%%%%%%%%%%%%%%%%%%%%%%%%%%%%%%%%%%%%%%%%%%%%%%%%%%%%%%%%%
\subsection{Interpreting the Slope}
\label{subsec:interpreting-regression-slope}
%%%%%%%%%%%%%%%%%%%%%%%%%%%%%%%%%%%%%%%%%%%%%%%%%%%%%%%%%%%%

The slope

\[
\beta_1
\]

represents the change in the conditional mean response associated with a
one-unit increase in \(X\).

Thus,

\[
\boxed{
\Delta
\mathbb{E}[Y\mid X]
=
\beta_1
}
\]

for a one-unit increase in \(X\).

If

\[
\beta_1>0,
\]

the mean response increases with \(X\).

If

\[
\beta_1<0,
\]

the mean response decreases with \(X\).

%%%%%%%%%%%%%%%%%%%%%%%%%%%%%%%%%%%%%%%%%%%%%%%%%%%%%%%%%%%%
\subsection{Interpreting the Intercept}
\label{subsec:interpreting-regression-intercept}
%%%%%%%%%%%%%%%%%%%%%%%%%%%%%%%%%%%%%%%%%%%%%%%%%%%%%%%%%%%%

The intercept

\[
\beta_0
\]

is the predicted mean response when

\[
X=0.
\]

Thus,

\[
\boxed{
\beta_0
=
\mathbb{E}[Y\mid X=0].
}
\]

However, the intercept may not have a meaningful scientific
interpretation if

\[
X=0
\]

is impossible or far outside the observed range.

%%%%%%%%%%%%%%%%%%%%%%%%%%%%%%%%%%%%%%%%%%%%%%%%%%%%%%%%%%%%
\subsection{Observed Data}
\label{subsec:observed-regression-data}
%%%%%%%%%%%%%%%%%%%%%%%%%%%%%%%%%%%%%%%%%%%%%%%%%%%%%%%%%%%%

Suppose the observed data are

\[
(x_1,y_1),
\ldots,
(x_n,y_n).
\]

The estimated regression line is

\[
\boxed{
\widehat{y}
=
b_0+b_1x.
}
\]

The coefficients

\[
b_0
\]

and

\[
b_1
\]

estimate

\[
\beta_0
\]

and

\[
\beta_1.
\]

%%%%%%%%%%%%%%%%%%%%%%%%%%%%%%%%%%%%%%%%%%%%%%%%%%%%%%%%%%%%
\subsection{Residuals}
\label{subsec:regression-residuals}
%%%%%%%%%%%%%%%%%%%%%%%%%%%%%%%%%%%%%%%%%%%%%%%%%%%%%%%%%%%%

For observation \(i\), the fitted value is

\[
\boxed{
\widehat{y}_i
=
b_0+b_1x_i.
}
\]

The residual is

\[
\boxed{
e_i
=
y_i-\widehat{y}_i.
}
\]

Thus residuals estimate the unobserved errors.

%%%%%%%%%%%%%%%%%%%%%%%%%%%%%%%%%%%%%%%%%%%%%%%%%%%%%%%%%%%%
\subsection{Ordinary Least Squares}
\label{subsec:ordinary-least-squares}
%%%%%%%%%%%%%%%%%%%%%%%%%%%%%%%%%%%%%%%%%%%%%%%%%%%%%%%%%%%%

Ordinary least squares, abbreviated OLS, chooses coefficients minimizing

\[
\boxed{
SSE
=
\sum_{i=1}^{n}
e_i^2.
}
\]

Thus,

\[
\boxed{
(b_0,b_1)
=
\arg\min_{a,b}
\sum_{i=1}^{n}
(y_i-a-bx_i)^2.
}
\]

This is called the least-squares criterion.

%%%%%%%%%%%%%%%%%%%%%%%%%%%%%%%%%%%%%%%%%%%%%%%%%%%%%%%%%%%%
\subsection{Deriving the Least-Squares Equations}
\label{subsec:derive-least-squares}
%%%%%%%%%%%%%%%%%%%%%%%%%%%%%%%%%%%%%%%%%%%%%%%%%%%%%%%%%%%%

Define

\[
Q(a,b)
=
\sum_{i=1}^{n}
(y_i-a-bx_i)^2.
\]

Differentiate with respect to \(a\):

\[
\frac{\partial Q}{\partial a}
=
-2
\sum_{i=1}^{n}
(y_i-a-bx_i).
\]

Set equal to zero:

\[
\sum y_i
=
na+b\sum x_i.
\]

Differentiate with respect to \(b\):

\[
\frac{\partial Q}{\partial b}
=
-2
\sum
x_i(y_i-a-bx_i).
\]

Set equal to zero:

\[
\sum x_iy_i
=
a\sum x_i
+
b\sum x_i^2.
\]

These are the normal equations.

%%%%%%%%%%%%%%%%%%%%%%%%%%%%%%%%%%%%%%%%%%%%%%%%%%%%%%%%%%%%
\subsection{Slope Estimator}
\label{subsec:slope-estimator}
%%%%%%%%%%%%%%%%%%%%%%%%%%%%%%%%%%%%%%%%%%%%%%%%%%%%%%%%%%%%

Solving the normal equations gives

\[
\boxed{
b_1
=
\frac{
\sum_{i=1}^{n}
(x_i-\overline{x})
(y_i-\overline{y})
}{
\sum_{i=1}^{n}
(x_i-\overline{x})^2
}.
}
\]

Define

\[
S_{xx}
=
\sum
(x_i-\overline{x})^2,
\]

and

\[
S_{xy}
=
\sum
(x_i-\overline{x})
(y_i-\overline{y}).
\]

Then

\[
\boxed{
b_1
=
\frac{S_{xy}}{S_{xx}}.
}
\]

%%%%%%%%%%%%%%%%%%%%%%%%%%%%%%%%%%%%%%%%%%%%%%%%%%%%%%%%%%%%
\subsection{Intercept Estimator}
\label{subsec:intercept-estimator}
%%%%%%%%%%%%%%%%%%%%%%%%%%%%%%%%%%%%%%%%%%%%%%%%%%%%%%%%%%%%

Once the slope is known,

\[
\boxed{
b_0
=
\overline{y}
-
b_1\overline{x}.
}
\]

Therefore the estimated regression line always passes through

\[
\boxed{
(\overline{x},\overline{y}).
}
\]

%%%%%%%%%%%%%%%%%%%%%%%%%%%%%%%%%%%%%%%%%%%%%%%%%%%%%%%%%%%%
\subsection{Worked Example: Least-Squares Line}
\label{subsec:worked-least-squares-line}
%%%%%%%%%%%%%%%%%%%%%%%%%%%%%%%%%%%%%%%%%%%%%%%%%%%%%%%%%%%%

\begin{workedexamplebox}{Fitting a Simple Regression}

Suppose

\[
x=(1,2,3)
\]

and

\[
y=(2,3,5).
\]

Then

\[
\overline{x}=2,
\qquad
\overline{y}=\frac{10}{3}.
\]

Compute

\[
S_{xx}
=
(1-2)^2
+
(2-2)^2
+
(3-2)^2
=
2.
\]

Also,

\[
S_{xy}
=
(1-2)
\left(
2-\frac{10}{3}
\right)
+
(2-2)
\left(
3-\frac{10}{3}
\right)
+
(3-2)
\left(
5-\frac{10}{3}
\right).
\]

Thus,

\[
S_{xy}
=
\frac43+\frac53
=
3.
\]

Therefore,

\[
b_1
=
\frac32.
\]

The intercept is

\[
b_0
=
\frac{10}{3}
-
\frac32(2)
=
\frac13.
\]

\begin{answerbox}
\[
\boxed{
\widehat{y}
=
\frac13
+
\frac32x.
}
\]
\end{answerbox}

\end{workedexamplebox}

%%%%%%%%%%%%%%%%%%%%%%%%%%%%%%%%%%%%%%%%%%%%%%%%%%%%%%%%%%%%
\subsection{Properties of OLS Residuals}
\label{subsec:ols-residual-properties}
%%%%%%%%%%%%%%%%%%%%%%%%%%%%%%%%%%%%%%%%%%%%%%%%%%%%%%%%%%%%

For simple regression with an intercept,

\[
\boxed{
\sum_{i=1}^{n}e_i=0.
}
\]

Also,

\[
\boxed{
\sum_{i=1}^{n}x_ie_i=0.
}
\]

And since fitted values are linear combinations of \(1\) and \(x_i\),

\[
\boxed{
\sum_{i=1}^{n}\widehat{y}_ie_i=0.
}
\]

Thus residuals are orthogonal to the fitted regression space.

%%%%%%%%%%%%%%%%%%%%%%%%%%%%%%%%%%%%%%%%%%%%%%%%%%%%%%%%%%%%
\subsection{Mean of the Fitted Values}
\label{subsec:mean-fitted-values}
%%%%%%%%%%%%%%%%%%%%%%%%%%%%%%%%%%%%%%%%%%%%%%%%%%%%%%%%%%%%

Since

\[
\sum e_i=0,
\]

we have

\[
\sum y_i
=
\sum\widehat{y}_i.
\]

Therefore,

\[
\boxed{
\overline{\widehat{y}}
=
\overline{y}.
}
\]

The fitted values and observed responses have the same sample mean when
an intercept is included.

%%%%%%%%%%%%%%%%%%%%%%%%%%%%%%%%%%%%%%%%%%%%%%%%%%%%%%%%%%%%
\subsection{Total Sum of Squares}
\label{subsec:total-sum-squares}
%%%%%%%%%%%%%%%%%%%%%%%%%%%%%%%%%%%%%%%%%%%%%%%%%%%%%%%%%%%%

The total variation in the response is

\[
\boxed{
SST
=
\sum_{i=1}^{n}
(y_i-\overline{y})^2.
}
\]

This measures variability of the observations around their mean.

%%%%%%%%%%%%%%%%%%%%%%%%%%%%%%%%%%%%%%%%%%%%%%%%%%%%%%%%%%%%
\subsection{Residual Sum of Squares}
\label{subsec:residual-sum-squares}
%%%%%%%%%%%%%%%%%%%%%%%%%%%%%%%%%%%%%%%%%%%%%%%%%%%%%%%%%%%%

The residual sum of squares is

\[
\boxed{
SSE
=
\sum_{i=1}^{n}
(y_i-\widehat{y}_i)^2.
}
\]

This is the amount of variability not explained by the fitted regression
line.

%%%%%%%%%%%%%%%%%%%%%%%%%%%%%%%%%%%%%%%%%%%%%%%%%%%%%%%%%%%%
\subsection{Regression Sum of Squares}
\label{subsec:regression-sum-squares}
%%%%%%%%%%%%%%%%%%%%%%%%%%%%%%%%%%%%%%%%%%%%%%%%%%%%%%%%%%%%

The regression sum of squares is

\[
\boxed{
SSR
=
\sum_{i=1}^{n}
(\widehat{y}_i-\overline{y})^2.
}
\]

This measures the variability of the fitted values around the response
mean.

%%%%%%%%%%%%%%%%%%%%%%%%%%%%%%%%%%%%%%%%%%%%%%%%%%%%%%%%%%%%
\subsection{ANOVA Decomposition}
\label{subsec:regression-anova-decomposition}
%%%%%%%%%%%%%%%%%%%%%%%%%%%%%%%%%%%%%%%%%%%%%%%%%%%%%%%%%%%%

For ordinary least squares with an intercept,

\[
\boxed{
SST
=
SSR
+
SSE.
}
\]

Thus,

\[
\boxed{
\text{total variation}
=
\text{explained variation}
+
\text{unexplained variation}.
}
\]

This identity forms the basis of regression ANOVA.

%%%%%%%%%%%%%%%%%%%%%%%%%%%%%%%%%%%%%%%%%%%%%%%%%%%%%%%%%%%%
\subsection{Coefficient of Determination}
\label{subsec:coefficient-determination}
%%%%%%%%%%%%%%%%%%%%%%%%%%%%%%%%%%%%%%%%%%%%%%%%%%%%%%%%%%%%

The coefficient of determination is

\[
\boxed{
R^2
=
\frac{SSR}{SST}.
}
\]

Equivalently,

\[
\boxed{
R^2
=
1-\frac{SSE}{SST}.
}
\]

When an intercept is included,

\[
\boxed{
0\leq R^2\leq1.
}
\]

%%%%%%%%%%%%%%%%%%%%%%%%%%%%%%%%%%%%%%%%%%%%%%%%%%%%%%%%%%%%
\subsection{Interpretation of \(R^2\)}
\label{subsec:r2-interpretation}
%%%%%%%%%%%%%%%%%%%%%%%%%%%%%%%%%%%%%%%%%%%%%%%%%%%%%%%%%%%%

An \(R^2\) value of

\[
0.80
\]

means that \(80\%\) of the observed sample variation around
\(\overline{y}\) is accounted for by the fitted regression model in the
sum-of-squares sense.

It does not mean that:

\[
\boxed{
80\%
\text{ of }Y
\text{ is caused by }X.
}
\]

Nor does it guarantee accurate prediction outside the observed sample.

%%%%%%%%%%%%%%%%%%%%%%%%%%%%%%%%%%%%%%%%%%%%%%%%%%%%%%%%%%%%
\subsection{Correlation and Simple Regression}
\label{subsec:correlation-simple-regression}
%%%%%%%%%%%%%%%%%%%%%%%%%%%%%%%%%%%%%%%%%%%%%%%%%%%%%%%%%%%%

Recall the sample correlation

\[
r
=
\frac{
S_{xy}
}{
\sqrt{
S_{xx}S_{yy}
}
},
\]

where

\[
S_{yy}
=
\sum
(y_i-\overline{y})^2.
\]

The simple-regression slope is

\[
b_1
=
\frac{S_{xy}}{S_{xx}}.
\]

Therefore,

\[
\boxed{
b_1
=
r
\frac{s_y}{s_x}.
}
\]

%%%%%%%%%%%%%%%%%%%%%%%%%%%%%%%%%%%%%%%%%%%%%%%%%%%%%%%%%%%%
\subsection{\(R^2\) and Correlation}
\label{subsec:r2-correlation}
%%%%%%%%%%%%%%%%%%%%%%%%%%%%%%%%%%%%%%%%%%%%%%%%%%%%%%%%%%%%

For simple linear regression with an intercept,

\[
\boxed{
R^2=r^2.
}
\]

Thus the squared Pearson correlation equals the coefficient of
determination.

This identity is specific to simple linear regression with an intercept.

%%%%%%%%%%%%%%%%%%%%%%%%%%%%%%%%%%%%%%%%%%%%%%%%%%%%%%%%%%%%
\subsection{Worked Example: Slope from Correlation}
\label{subsec:worked-slope-correlation}
%%%%%%%%%%%%%%%%%%%%%%%%%%%%%%%%%%%%%%%%%%%%%%%%%%%%%%%%%%%%

\begin{workedexamplebox}{Using Standard Deviations and Correlation}

Suppose

\[
r=0.75,
\qquad
s_x=4,
\qquad
s_y=10.
\]

Then

\[
b_1
=
r
\frac{s_y}{s_x}.
\]

Thus,

\begin{answerbox}
\[
\boxed{
b_1
=
0.75
\frac{10}{4}
=
1.875.
}
\]
\end{answerbox}

\end{workedexamplebox}

%%%%%%%%%%%%%%%%%%%%%%%%%%%%%%%%%%%%%%%%%%%%%%%%%%%%%%%%%%%%
\subsection{Classical Linear Regression Assumptions}
\label{subsec:classical-linear-regression-assumptions}
%%%%%%%%%%%%%%%%%%%%%%%%%%%%%%%%%%%%%%%%%%%%%%%%%%%%%%%%%%%%

A common simple linear regression model assumes

\[
\boxed{
Y_i
=
\beta_0+\beta_1x_i+\varepsilon_i,
}
\]

with

\[
\boxed{
\mathbb{E}[\varepsilon_i]=0,
}
\]

\[
\boxed{
\operatorname{Var}(\varepsilon_i)=\sigma^2,
}
\]

and

\[
\boxed{
\operatorname{Cov}(\varepsilon_i,\varepsilon_j)=0
\quad
(i\neq j).
}
\]

For exact small-sample \(t\) and \(F\) inference, one commonly also
assumes

\[
\boxed{
\varepsilon_i
\overset{\text{ind.}}{\sim}
N(0,\sigma^2).
}
\]

%%%%%%%%%%%%%%%%%%%%%%%%%%%%%%%%%%%%%%%%%%%%%%%%%%%%%%%%%%%%
\subsection{Linearity Assumption}
\label{subsec:regression-linearity-assumption}
%%%%%%%%%%%%%%%%%%%%%%%%%%%%%%%%%%%%%%%%%%%%%%%%%%%%%%%%%%%%

The simple linear model assumes

\[
\boxed{
\mathbb{E}[Y\mid X=x]
=
\beta_0+\beta_1x.
}
\]

This means linearity in the coefficients and in the conditional mean
structure being modeled.

A visibly curved relationship may require transformations, polynomial
terms, splines, or another model.

%%%%%%%%%%%%%%%%%%%%%%%%%%%%%%%%%%%%%%%%%%%%%%%%%%%%%%%%%%%%
\subsection{Zero Conditional Mean}
\label{subsec:zero-conditional-mean}
%%%%%%%%%%%%%%%%%%%%%%%%%%%%%%%%%%%%%%%%%%%%%%%%%%%%%%%%%%%%

A central assumption is

\[
\boxed{
\mathbb{E}[\varepsilon\mid X]=0.
}
\]

This says that once \(X\) is known, the remaining error has no systematic
mean pattern.

Violation of this condition can arise from:

\[
\boxed{
\text{omitted variables},
}
\]

\[
\boxed{
\text{simultaneity},
}
\]

or

\[
\boxed{
\text{selection effects}.
}
\]

%%%%%%%%%%%%%%%%%%%%%%%%%%%%%%%%%%%%%%%%%%%%%%%%%%%%%%%%%%%%
\subsection{Homoskedasticity}
\label{subsec:homoskedasticity}
%%%%%%%%%%%%%%%%%%%%%%%%%%%%%%%%%%%%%%%%%%%%%%%%%%%%%%%%%%%%

Homoskedasticity means

\[
\boxed{
\operatorname{Var}
(
\varepsilon_i\mid X_i
)
=
\sigma^2
}
\]

for all observations.

Thus the conditional error variance is constant.

If instead the variance changes with \(X\), the errors are
heteroskedastic.

%%%%%%%%%%%%%%%%%%%%%%%%%%%%%%%%%%%%%%%%%%%%%%%%%%%%%%%%%%%%
\subsection{Independence of Errors}
\label{subsec:independence-errors-regression}
%%%%%%%%%%%%%%%%%%%%%%%%%%%%%%%%%%%%%%%%%%%%%%%%%%%%%%%%%%%%

In many regression settings one assumes observations are independent.

This may fail for:

\[
\boxed{
\text{time series},
}
\]

\[
\boxed{
\text{clustered observations},
}
\]

\[
\boxed{
\text{spatial data},
}
\]

or

\[
\boxed{
\text{repeated measurements}.
}
\]

Dependence affects standard errors and inference.

%%%%%%%%%%%%%%%%%%%%%%%%%%%%%%%%%%%%%%%%%%%%%%%%%%%%%%%%%%%%
\subsection{Normality of Errors}
\label{subsec:normality-errors-regression}
%%%%%%%%%%%%%%%%%%%%%%%%%%%%%%%%%%%%%%%%%%%%%%%%%%%%%%%%%%%%

Normality is not required to compute least-squares estimates.

It is used to obtain exact finite-sample distributions for many classical
tests and intervals.

For large samples, asymptotic methods may remain useful under weaker
conditions.

Thus:

\[
\boxed{
\text{normality is mainly an inference assumption, not an OLS existence
requirement}.
}
\]

%%%%%%%%%%%%%%%%%%%%%%%%%%%%%%%%%%%%%%%%%%%%%%%%%%%%%%%%%%%%
\subsection{Expected Value of the OLS Slope}
\label{subsec:expected-ols-slope}
%%%%%%%%%%%%%%%%%%%%%%%%%%%%%%%%%%%%%%%%%%%%%%%%%%%%%%%%%%%%

For fixed \(x_i\),

\[
b_1
=
\beta_1
+
\frac{
\sum
(x_i-\overline{x})\varepsilon_i
}{
S_{xx}
}.
\]

If

\[
\mathbb{E}[\varepsilon_i]=0,
\]

then

\[
\boxed{
\mathbb{E}[b_1]
=
\beta_1.
}
\]

Thus the OLS slope is unbiased under the standard zero-mean assumptions.

%%%%%%%%%%%%%%%%%%%%%%%%%%%%%%%%%%%%%%%%%%%%%%%%%%%%%%%%%%%%
\subsection{Variance of the OLS Slope}
\label{subsec:variance-ols-slope}
%%%%%%%%%%%%%%%%%%%%%%%%%%%%%%%%%%%%%%%%%%%%%%%%%%%%%%%%%%%%

Under independent homoskedastic errors,

\[
\boxed{
\operatorname{Var}(b_1)
=
\frac{
\sigma^2
}{
S_{xx}
}.
}
\]

Therefore,

\[
\boxed{
\operatorname{SE}(b_1)
=
\frac{
\sigma
}{
\sqrt{S_{xx}}
}.
}
\]

Greater spread in the explanatory values reduces slope uncertainty.

%%%%%%%%%%%%%%%%%%%%%%%%%%%%%%%%%%%%%%%%%%%%%%%%%%%%%%%%%%%%
\subsection{Variance of the OLS Intercept}
\label{subsec:variance-ols-intercept}
%%%%%%%%%%%%%%%%%%%%%%%%%%%%%%%%%%%%%%%%%%%%%%%%%%%%%%%%%%%%

Under the standard assumptions,

\[
\boxed{
\operatorname{Var}(b_0)
=
\sigma^2
\left[
\frac1n
+
\frac{
\overline{x}^2
}{
S_{xx}
}
\right].
}
\]

Thus the intercept may be estimated imprecisely when \(x=0\) lies far
from the center of the observed \(x\)-values.

%%%%%%%%%%%%%%%%%%%%%%%%%%%%%%%%%%%%%%%%%%%%%%%%%%%%%%%%%%%%
\subsection{Covariance Between Intercept and Slope}
\label{subsec:covariance-intercept-slope}
%%%%%%%%%%%%%%%%%%%%%%%%%%%%%%%%%%%%%%%%%%%%%%%%%%%%%%%%%%%%

The covariance is

\[
\boxed{
\operatorname{Cov}(b_0,b_1)
=
-
\frac{
\overline{x}\sigma^2
}{
S_{xx}
}.
}
\]

If \(X\) is centered so that

\[
\overline{x}=0,
\]

then

\[
\boxed{
\operatorname{Cov}(b_0,b_1)=0.
}
\]

Centering can therefore improve coefficient interpretation and numerical
behavior.

%%%%%%%%%%%%%%%%%%%%%%%%%%%%%%%%%%%%%%%%%%%%%%%%%%%%%%%%%%%%
\subsection{Residual Variance Estimator}
\label{subsec:residual-variance-estimator}
%%%%%%%%%%%%%%%%%%%%%%%%%%%%%%%%%%%%%%%%%%%%%%%%%%%%%%%%%%%%

In simple regression, two parameters

\[
\beta_0,
\qquad
\beta_1
\]

are estimated.

Therefore the residual degrees of freedom are

\[
\boxed{
n-2.
}
\]

The error variance is estimated by

\[
\boxed{
S^2
=
\frac{
SSE
}{
n-2
}.
}
\]

The residual standard error is

\[
\boxed{
S
=
\sqrt{
\frac{
SSE
}{
n-2
}
}.
}
\]

%%%%%%%%%%%%%%%%%%%%%%%%%%%%%%%%%%%%%%%%%%%%%%%%%%%%%%%%%%%%
\subsection{Estimated Standard Error of the Slope}
\label{subsec:estimated-se-slope}
%%%%%%%%%%%%%%%%%%%%%%%%%%%%%%%%%%%%%%%%%%%%%%%%%%%%%%%%%%%%

Replace \(\sigma\) by \(S\):

\[
\boxed{
SE(b_1)
=
\frac{
S
}{
\sqrt{S_{xx}}
}.
}
\]

Similarly,

\[
\boxed{
SE(b_0)
=
S
\sqrt{
\frac1n
+
\frac{
\overline{x}^2
}{
S_{xx}
}
}.
}
\]

These are used in coefficient confidence intervals and tests.

%%%%%%%%%%%%%%%%%%%%%%%%%%%%%%%%%%%%%%%%%%%%%%%%%%%%%%%%%%%%
\subsection{Sampling Distribution of the Slope}
\label{subsec:sampling-distribution-slope}
%%%%%%%%%%%%%%%%%%%%%%%%%%%%%%%%%%%%%%%%%%%%%%%%%%%%%%%%%%%%

Under independent normal errors,

\[
b_1
\]

is normally distributed:

\[
\boxed{
b_1
\sim
N
\left(
\beta_1,
\frac{
\sigma^2
}{
S_{xx}
}
\right).
}
\]

After replacing \(\sigma\) with \(S\),

\[
\boxed{
\frac{
b_1-\beta_1
}{
SE(b_1)
}
\sim
t_{n-2}.
}
\]

%%%%%%%%%%%%%%%%%%%%%%%%%%%%%%%%%%%%%%%%%%%%%%%%%%%%%%%%%%%%
\subsection{Confidence Interval for the Slope}
\label{subsec:confidence-interval-slope}
%%%%%%%%%%%%%%%%%%%%%%%%%%%%%%%%%%%%%%%%%%%%%%%%%%%%%%%%%%%%

A \(100(1-\alpha)\%\) confidence interval for \(\beta_1\) is

\[
\boxed{
b_1
\pm
t^*_{n-2}
SE(b_1).
}
\]

This interval quantifies uncertainty in the estimated change in mean
response per unit change in \(X\).

%%%%%%%%%%%%%%%%%%%%%%%%%%%%%%%%%%%%%%%%%%%%%%%%%%%%%%%%%%%%
\subsection{Hypothesis Test for the Slope}
\label{subsec:hypothesis-test-slope}
%%%%%%%%%%%%%%%%%%%%%%%%%%%%%%%%%%%%%%%%%%%%%%%%%%%%%%%%%%%%

A common hypothesis is

\[
\boxed{
H_0:\beta_1=0.
}
\]

Against the two-sided alternative

\[
\boxed{
H_A:\beta_1\neq0.
}
\]

The test statistic is

\[
\boxed{
T
=
\frac{
b_1
}{
SE(b_1)
}.
}
\]

Under \(H_0\) and the normal-error model,

\[
\boxed{
T\sim t_{n-2}.
}
\]

%%%%%%%%%%%%%%%%%%%%%%%%%%%%%%%%%%%%%%%%%%%%%%%%%%%%%%%%%%%%
\subsection{Meaning of \(H_0:\beta_1=0\)}
\label{subsec:meaning-zero-slope}
%%%%%%%%%%%%%%%%%%%%%%%%%%%%%%%%%%%%%%%%%%%%%%%%%%%%%%%%%%%%

In simple linear regression,

\[
\beta_1=0
\]

means

\[
\boxed{
\mathbb{E}[Y\mid X=x]
=
\beta_0
}
\]

for all \(x\) under the model.

Thus the conditional mean does not change linearly with \(X\).

This is a statement about linear association, not necessarily complete
statistical independence.

%%%%%%%%%%%%%%%%%%%%%%%%%%%%%%%%%%%%%%%%%%%%%%%%%%%%%%%%%%%%
\subsection{Worked Example: Testing the Slope}
\label{subsec:worked-testing-slope}
%%%%%%%%%%%%%%%%%%%%%%%%%%%%%%%%%%%%%%%%%%%%%%%%%%%%%%%%%%%%

\begin{workedexamplebox}{Regression Coefficient Test}

Suppose

\[
b_1=2.4
\]

and

\[
SE(b_1)=0.8.
\]

Test

\[
H_0:\beta_1=0.
\]

Then

\[
t
=
\frac{2.4}{0.8}
=
3.
\]

\begin{answerbox}
\[
\boxed{
t=3.
}
\]
\end{answerbox}

The degrees of freedom are \(n-2\).

\end{workedexamplebox}

%%%%%%%%%%%%%%%%%%%%%%%%%%%%%%%%%%%%%%%%%%%%%%%%%%%%%%%%%%%%
\subsection{Confidence Interval for the Mean Response}
\label{subsec:ci-mean-response}
%%%%%%%%%%%%%%%%%%%%%%%%%%%%%%%%%%%%%%%%%%%%%%%%%%%%%%%%%%%%

At a specified predictor value

\[
x_0,
\]

the estimated mean response is

\[
\boxed{
\widehat{y}_0
=
b_0+b_1x_0.
}
\]

Its estimated standard error is

\[
\boxed{
SE_{\text{mean}}
=
S
\sqrt{
\frac1n
+
\frac{
(x_0-\overline{x})^2
}{
S_{xx}
}
}.
}
\]

Therefore,

\[
\boxed{
\widehat{y}_0
\pm
t^*_{n-2}
SE_{\text{mean}}
}
\]

is a confidence interval for the mean response at \(x_0\).

%%%%%%%%%%%%%%%%%%%%%%%%%%%%%%%%%%%%%%%%%%%%%%%%%%%%%%%%%%%%
\subsection{Prediction Interval for a New Observation}
\label{subsec:prediction-interval-regression}
%%%%%%%%%%%%%%%%%%%%%%%%%%%%%%%%%%%%%%%%%%%%%%%%%%%%%%%%%%%%

A future observation at \(x_0\) contains additional individual error.

Thus its prediction standard error is

\[
\boxed{
SE_{\text{pred}}
=
S
\sqrt{
1
+
\frac1n
+
\frac{
(x_0-\overline{x})^2
}{
S_{xx}
}
}.
}
\]

The prediction interval is

\[
\boxed{
\widehat{y}_0
\pm
t^*_{n-2}
SE_{\text{pred}}.
}
\]

%%%%%%%%%%%%%%%%%%%%%%%%%%%%%%%%%%%%%%%%%%%%%%%%%%%%%%%%%%%%
\subsection{Confidence Versus Prediction Interval}
\label{subsec:confidence-vs-prediction-regression}
%%%%%%%%%%%%%%%%%%%%%%%%%%%%%%%%%%%%%%%%%%%%%%%%%%%%%%%%%%%%

The mean-response confidence interval estimates

\[
\mathbb{E}[Y\mid X=x_0].
\]

The prediction interval targets an individual future value

\[
Y_{\text{new}}.
\]

Because

\[
SE_{\text{pred}}
>
SE_{\text{mean}},
\]

the prediction interval is wider.

The extra \(1\) inside the square root represents individual-response
variability.

%%%%%%%%%%%%%%%%%%%%%%%%%%%%%%%%%%%%%%%%%%%%%%%%%%%%%%%%%%%%
\subsection{Where Prediction Is Most Precise}
\label{subsec:prediction-precision-location}
%%%%%%%%%%%%%%%%%%%%%%%%%%%%%%%%%%%%%%%%%%%%%%%%%%%%%%%%%%%%

Both interval formulas contain

\[
(x_0-\overline{x})^2.
\]

Therefore uncertainty is smallest near

\[
\boxed{
x_0=\overline{x}.
}
\]

Uncertainty increases as \(x_0\) moves farther from the center of the
observed predictor values.

%%%%%%%%%%%%%%%%%%%%%%%%%%%%%%%%%%%%%%%%%%%%%%%%%%%%%%%%%%%%
\subsection{Extrapolation}
\label{subsec:regression-extrapolation}
%%%%%%%%%%%%%%%%%%%%%%%%%%%%%%%%%%%%%%%%%%%%%%%%%%%%%%%%%%%%

Extrapolation uses a regression model to predict outside the observed
range of explanatory variables.

For example, if observed \(X\)-values range from

\[
10
\]

to

\[
20,
\]

predicting at

\[
X=100
\]

is extrapolation.

The assumed relationship may not continue outside the observed range.

Therefore:

\[
\boxed{
\text{extrapolation can be highly unreliable}.
}
\]

%%%%%%%%%%%%%%%%%%%%%%%%%%%%%%%%%%%%%%%%%%%%%%%%%%%%%%%%%%%%
\subsection{Multiple Linear Regression}
\label{subsec:multiple-linear-regression}
%%%%%%%%%%%%%%%%%%%%%%%%%%%%%%%%%%%%%%%%%%%%%%%%%%%%%%%%%%%%

With several explanatory variables,

\[
\boxed{
Y_i
=
\beta_0
+
\beta_1X_{i1}
+
\cdots
+
\beta_pX_{ip}
+
\varepsilon_i.
}
\]

The conditional mean is

\[
\boxed{
\mathbb{E}
[
Y_i
\mid
X_{i1},\ldots,X_{ip}
]
=
\beta_0
+
\sum_{j=1}^{p}
\beta_jX_{ij}.
}
\]

%%%%%%%%%%%%%%%%%%%%%%%%%%%%%%%%%%%%%%%%%%%%%%%%%%%%%%%%%%%%
\subsection{Interpretation of a Partial Regression Coefficient}
\label{subsec:partial-regression-coefficient}
%%%%%%%%%%%%%%%%%%%%%%%%%%%%%%%%%%%%%%%%%%%%%%%%%%%%%%%%%%%%

In multiple regression,

\[
\beta_j
\]

represents the change in conditional mean response associated with a
one-unit increase in \(X_j\), while holding the other included
explanatory variables fixed.

Thus:

\[
\boxed{
\beta_j
=
\text{partial association of }X_j\text{ with }Y
\text{ given the other predictors}.
}
\]

This interpretation depends on the model specification.

%%%%%%%%%%%%%%%%%%%%%%%%%%%%%%%%%%%%%%%%%%%%%%%%%%%%%%%%%%%%
\subsection{Matrix Form of Regression}
\label{subsec:matrix-form-regression}
%%%%%%%%%%%%%%%%%%%%%%%%%%%%%%%%%%%%%%%%%%%%%%%%%%%%%%%%%%%%

Multiple regression can be written compactly as

\[
\boxed{
\mathbf{y}
=
X\boldsymbol{\beta}
+
\boldsymbol{\varepsilon}.
}
\]

Here,

\[
\mathbf{y}
=
\begin{pmatrix}
y_1\\
\vdots\\
y_n
\end{pmatrix},
\]

\[
\boldsymbol{\beta}
=
\begin{pmatrix}
\beta_0\\
\beta_1\\
\vdots\\
\beta_p
\end{pmatrix},
\]

and \(X\) is the design matrix.

%%%%%%%%%%%%%%%%%%%%%%%%%%%%%%%%%%%%%%%%%%%%%%%%%%%%%%%%%%%%
\subsection{Design Matrix}
\label{subsec:design-matrix-regression}
%%%%%%%%%%%%%%%%%%%%%%%%%%%%%%%%%%%%%%%%%%%%%%%%%%%%%%%%%%%%

If the model contains an intercept, the design matrix is

\[
\boxed{
X
=
\begin{pmatrix}
1 & x_{11} & \cdots & x_{1p}\\
1 & x_{21} & \cdots & x_{2p}\\
\vdots & \vdots & \ddots & \vdots\\
1 & x_{n1} & \cdots & x_{np}
\end{pmatrix}.
}
\]

Thus,

\[
X
\]

has

\[
n
\]

rows and

\[
p+1
\]

columns.

%%%%%%%%%%%%%%%%%%%%%%%%%%%%%%%%%%%%%%%%%%%%%%%%%%%%%%%%%%%%
\subsection{Least Squares in Matrix Form}
\label{subsec:least-squares-matrix}
%%%%%%%%%%%%%%%%%%%%%%%%%%%%%%%%%%%%%%%%%%%%%%%%%%%%%%%%%%%%

OLS minimizes

\[
\boxed{
\|\mathbf{y}-X\boldsymbol{\beta}\|_2^2.
}
\]

Differentiate with respect to \(\boldsymbol{\beta}\):

\[
-2X^T
(
\mathbf{y}
-
X\boldsymbol{\beta}
)
=
0.
\]

Therefore the normal equations are

\[
\boxed{
X^TX
\widehat{\boldsymbol{\beta}}
=
X^T\mathbf{y}.
}
\]

%%%%%%%%%%%%%%%%%%%%%%%%%%%%%%%%%%%%%%%%%%%%%%%%%%%%%%%%%%%%
\subsection{OLS Matrix Estimator}
\label{subsec:ols-matrix-estimator}
%%%%%%%%%%%%%%%%%%%%%%%%%%%%%%%%%%%%%%%%%%%%%%%%%%%%%%%%%%%%

If

\[
X^TX
\]

is invertible, then

\[
\boxed{
\widehat{\boldsymbol{\beta}}
=
(X^TX)^{-1}
X^T\mathbf{y}.
}
\]

This is the ordinary least-squares estimator.

Invertibility requires the columns of \(X\) to be linearly independent.

%%%%%%%%%%%%%%%%%%%%%%%%%%%%%%%%%%%%%%%%%%%%%%%%%%%%%%%%%%%%
\subsection{Full Column Rank}
\label{subsec:full-column-rank-regression}
%%%%%%%%%%%%%%%%%%%%%%%%%%%%%%%%%%%%%%%%%%%%%%%%%%%%%%%%%%%%

If the design matrix has

\[
p+1
\]

columns, full column rank means

\[
\boxed{
\operatorname{rank}(X)
=
p+1.
}
\]

If this fails, some predictors are exact linear combinations of others.

Then the individual regression coefficients are not uniquely
identified by ordinary least squares.

%%%%%%%%%%%%%%%%%%%%%%%%%%%%%%%%%%%%%%%%%%%%%%%%%%%%%%%%%%%%
\subsection{Fitted Values in Matrix Form}
\label{subsec:fitted-values-matrix}
%%%%%%%%%%%%%%%%%%%%%%%%%%%%%%%%%%%%%%%%%%%%%%%%%%%%%%%%%%%%

The fitted vector is

\[
\boxed{
\widehat{\mathbf{y}}
=
X
\widehat{\boldsymbol{\beta}}.
}
\]

Substitute the OLS estimator:

\[
\widehat{\mathbf{y}}
=
X
(X^TX)^{-1}
X^T
\mathbf{y}.
\]

Define the hat matrix

\[
\boxed{
H
=
X
(X^TX)^{-1}
X^T.
}
\]

Then

\[
\boxed{
\widehat{\mathbf{y}}
=
H\mathbf{y}.
}
\]

%%%%%%%%%%%%%%%%%%%%%%%%%%%%%%%%%%%%%%%%%%%%%%%%%%%%%%%%%%%%
\subsection{Why It Is Called the Hat Matrix}
\label{subsec:hat-matrix-name}
%%%%%%%%%%%%%%%%%%%%%%%%%%%%%%%%%%%%%%%%%%%%%%%%%%%%%%%%%%%%

The matrix \(H\) transforms the response vector

\[
\mathbf{y}
\]

into the fitted vector

\[
\widehat{\mathbf{y}}.
\]

Thus it ``puts the hat on \(y\)'':

\[
\boxed{
\widehat{\mathbf{y}}
=
H\mathbf{y}.
}
\]

%%%%%%%%%%%%%%%%%%%%%%%%%%%%%%%%%%%%%%%%%%%%%%%%%%%%%%%%%%%%
\subsection{Projection Interpretation}
\label{subsec:projection-interpretation-regression}
%%%%%%%%%%%%%%%%%%%%%%%%%%%%%%%%%%%%%%%%%%%%%%%%%%%%%%%%%%%%

The fitted vector is the orthogonal projection of \(\mathbf{y}\) onto
the column space of \(X\).

Thus

\[
\boxed{
\widehat{\mathbf{y}}
\in
\operatorname{Col}(X).
}
\]

The residual vector

\[
\mathbf{e}
=
\mathbf{y}
-
\widehat{\mathbf{y}}
\]

is orthogonal to that column space:

\[
\boxed{
X^T\mathbf{e}
=
0.
}
\]

This is the geometric meaning of the normal equations.

%%%%%%%%%%%%%%%%%%%%%%%%%%%%%%%%%%%%%%%%%%%%%%%%%%%%%%%%%%%%
\subsection{Properties of the Hat Matrix}
\label{subsec:hat-matrix-properties}
%%%%%%%%%%%%%%%%%%%%%%%%%%%%%%%%%%%%%%%%%%%%%%%%%%%%%%%%%%%%

The hat matrix satisfies

\[
\boxed{
H^T=H
}
\]

and

\[
\boxed{
H^2=H.
}
\]

Thus \(H\) is symmetric and idempotent.

Its trace is

\[
\boxed{
\operatorname{tr}(H)
=
p+1
}
\]

when \(X\) has full column rank.

%%%%%%%%%%%%%%%%%%%%%%%%%%%%%%%%%%%%%%%%%%%%%%%%%%%%%%%%%%%%
\subsection{Residual Maker Matrix}
\label{subsec:residual-maker}
%%%%%%%%%%%%%%%%%%%%%%%%%%%%%%%%%%%%%%%%%%%%%%%%%%%%%%%%%%%%

Define

\[
\boxed{
M
=
I-H.
}
\]

Then

\[
\boxed{
\mathbf{e}
=
M\mathbf{y}.
}
\]

The matrix \(M\) is also symmetric and idempotent.

Its rank is

\[
\boxed{
n-p-1.
}
\]

This equals the residual degrees of freedom.

%%%%%%%%%%%%%%%%%%%%%%%%%%%%%%%%%%%%%%%%%%%%%%%%%%%%%%%%%%%%
\subsection{Unbiasedness in Multiple Regression}
\label{subsec:unbiasedness-multiple-regression}
%%%%%%%%%%%%%%%%%%%%%%%%%%%%%%%%%%%%%%%%%%%%%%%%%%%%%%%%%%%%

Assume

\[
\mathbf{y}
=
X\boldsymbol{\beta}
+
\boldsymbol{\varepsilon}
\]

with

\[
\mathbb{E}
[
\boldsymbol{\varepsilon}\mid X
]
=
\mathbf{0}.
\]

Then

\[
\widehat{\boldsymbol{\beta}}
=
\boldsymbol{\beta}
+
(X^TX)^{-1}
X^T
\boldsymbol{\varepsilon}.
\]

Therefore,

\[
\boxed{
\mathbb{E}
[
\widehat{\boldsymbol{\beta}}
\mid X
]
=
\boldsymbol{\beta}.
}
\]

Thus OLS is conditionally unbiased.

%%%%%%%%%%%%%%%%%%%%%%%%%%%%%%%%%%%%%%%%%%%%%%%%%%%%%%%%%%%%
\subsection{Covariance Matrix of OLS}
\label{subsec:covariance-matrix-ols}
%%%%%%%%%%%%%%%%%%%%%%%%%%%%%%%%%%%%%%%%%%%%%%%%%%%%%%%%%%%%

Under homoskedastic independent errors,

\[
\operatorname{Var}
(
\boldsymbol{\varepsilon}\mid X
)
=
\sigma^2I.
\]

Therefore,

\[
\boxed{
\operatorname{Var}
(
\widehat{\boldsymbol{\beta}}
\mid X
)
=
\sigma^2
(X^TX)^{-1}.
}
\]

This matrix contains the variances and covariances of the estimated
coefficients.

%%%%%%%%%%%%%%%%%%%%%%%%%%%%%%%%%%%%%%%%%%%%%%%%%%%%%%%%%%%%
\subsection{Residual Variance in Multiple Regression}
\label{subsec:residual-variance-multiple-regression}
%%%%%%%%%%%%%%%%%%%%%%%%%%%%%%%%%%%%%%%%%%%%%%%%%%%%%%%%%%%%

With \(p\) predictors and an intercept, the number of estimated
coefficients is

\[
p+1.
\]

Therefore the residual degrees of freedom are

\[
\boxed{
n-p-1.
}
\]

The error variance estimator is

\[
\boxed{
S^2
=
\frac{
SSE
}{
n-p-1
}.
}
\]

%%%%%%%%%%%%%%%%%%%%%%%%%%%%%%%%%%%%%%%%%%%%%%%%%%%%%%%%%%%%
\subsection{Coefficient Standard Errors}
\label{subsec:coefficient-standard-errors}
%%%%%%%%%%%%%%%%%%%%%%%%%%%%%%%%%%%%%%%%%%%%%%%%%%%%%%%%%%%%

Let

\[
C
=
(X^TX)^{-1}.
\]

Then

\[
\operatorname{Var}
(
\widehat{\beta}_j
\mid X
)
=
\sigma^2C_{jj}.
\]

The estimated standard error is

\[
\boxed{
SE(\widehat{\beta}_j)
=
S\sqrt{C_{jj}}.
}
\]

%%%%%%%%%%%%%%%%%%%%%%%%%%%%%%%%%%%%%%%%%%%%%%%%%%%%%%%%%%%%
\subsection{Coefficient \(t\)-Tests}
\label{subsec:coefficient-t-tests}
%%%%%%%%%%%%%%%%%%%%%%%%%%%%%%%%%%%%%%%%%%%%%%%%%%%%%%%%%%%%

To test

\[
H_0:
\beta_j=\beta_{j0},
\]

use

\[
\boxed{
T
=
\frac{
\widehat{\beta}_j-\beta_{j0}
}{
SE(\widehat{\beta}_j)
}.
}
\]

Under normal homoskedastic errors,

\[
\boxed{
T
\sim
t_{n-p-1}.
}
\]

The most common null is

\[
\beta_{j0}=0.
\]

%%%%%%%%%%%%%%%%%%%%%%%%%%%%%%%%%%%%%%%%%%%%%%%%%%%%%%%%%%%%
\subsection{Confidence Interval for a Regression Coefficient}
\label{subsec:ci-regression-coefficient}
%%%%%%%%%%%%%%%%%%%%%%%%%%%%%%%%%%%%%%%%%%%%%%%%%%%%%%%%%%%%

A \(100(1-\alpha)\%\) confidence interval is

\[
\boxed{
\widehat{\beta}_j
\pm
t^*_{n-p-1}
SE(\widehat{\beta}_j).
}
\]

This interval describes plausible values for the partial regression
coefficient under the model assumptions.

%%%%%%%%%%%%%%%%%%%%%%%%%%%%%%%%%%%%%%%%%%%%%%%%%%%%%%%%%%%%
\subsection{Overall \(F\)-Test}
\label{subsec:overall-f-test-regression}
%%%%%%%%%%%%%%%%%%%%%%%%%%%%%%%%%%%%%%%%%%%%%%%%%%%%%%%%%%%%

For a model with \(p\) predictors, a common overall hypothesis is

\[
\boxed{
H_0:
\beta_1
=
\cdots
=
\beta_p
=
0.
}
\]

The alternative is that at least one slope is nonzero.

Define

\[
MSR
=
\frac{
SSR
}{p}
\]

and

\[
MSE
=
\frac{
SSE
}{
n-p-1
}.
\]

Then

\[
\boxed{
F
=
\frac{
MSR
}{
MSE
}.
}
\]

Under the classical null,

\[
\boxed{
F
\sim
F_{p,n-p-1}.
}
\]

%%%%%%%%%%%%%%%%%%%%%%%%%%%%%%%%%%%%%%%%%%%%%%%%%%%%%%%%%%%%
\subsection{Regression ANOVA Table}
\label{subsec:regression-anova-table}
%%%%%%%%%%%%%%%%%%%%%%%%%%%%%%%%%%%%%%%%%%%%%%%%%%%%%%%%%%%%

A regression ANOVA table has the structure

\[
\begin{array}{c|c|c|c|c}
\text{Source}
&
\text{df}
&
\text{SS}
&
\text{MS}
&
F
\\
\hline
\text{Regression}
&
p
&
SSR
&
SSR/p
&
MSR/MSE
\\
\text{Error}
&
n-p-1
&
SSE
&
SSE/(n-p-1)
&
\\
\text{Total}
&
n-1
&
SST
&
&
\end{array}
\]

This directly connects regression with analysis of variance.

%%%%%%%%%%%%%%%%%%%%%%%%%%%%%%%%%%%%%%%%%%%%%%%%%%%%%%%%%%%%
\subsection{Adjusted \(R^2\)}
\label{subsec:adjusted-r-squared}
%%%%%%%%%%%%%%%%%%%%%%%%%%%%%%%%%%%%%%%%%%%%%%%%%%%%%%%%%%%%

Ordinary \(R^2\) never decreases when more predictors are added.

Adjusted \(R^2\) penalizes model complexity:

\[
\boxed{
R_{\mathrm{adj}}^2
=
1
-
\frac{
SSE/(n-p-1)
}{
SST/(n-1)
}.
}
\]

Equivalently,

\[
\boxed{
R_{\mathrm{adj}}^2
=
1
-
(1-R^2)
\frac{
n-1
}{
n-p-1
}.
}
\]

Adjusted \(R^2\) can decrease when an unhelpful variable is added.

%%%%%%%%%%%%%%%%%%%%%%%%%%%%%%%%%%%%%%%%%%%%%%%%%%%%%%%%%%%%
\subsection{Indicator Variables}
\label{subsec:indicator-variables-regression}
%%%%%%%%%%%%%%%%%%%%%%%%%%%%%%%%%%%%%%%%%%%%%%%%%%%%%%%%%%%%

A categorical variable with two levels can be represented using an
indicator variable.

For example,

\[
D
=
\begin{cases}
1,&\text{group A},\\
0,&\text{group B}.
\end{cases}
\]

Consider

\[
\boxed{
Y
=
\beta_0
+
\beta_1D
+
\varepsilon.
}
\]

Then for group B,

\[
\mathbb{E}[Y\mid D=0]
=
\beta_0.
\]

For group A,

\[
\mathbb{E}[Y\mid D=1]
=
\beta_0+\beta_1.
\]

Therefore,

\[
\boxed{
\beta_1
=
\text{difference in group means}.
}
\]

%%%%%%%%%%%%%%%%%%%%%%%%%%%%%%%%%%%%%%%%%%%%%%%%%%%%%%%%%%%%
\subsection{Categorical Predictors with Several Levels}
\label{subsec:categorical-predictors-regression}
%%%%%%%%%%%%%%%%%%%%%%%%%%%%%%%%%%%%%%%%%%%%%%%%%%%%%%%%%%%%

Suppose a categorical predictor has \(k\) levels.

With an intercept, one commonly uses

\[
\boxed{
k-1
}
\]

indicator variables.

One category serves as the reference category.

Each indicator coefficient measures a difference relative to that
reference level.

Including all \(k\) indicators together with an intercept creates exact
linear dependence.

%%%%%%%%%%%%%%%%%%%%%%%%%%%%%%%%%%%%%%%%%%%%%%%%%%%%%%%%%%%%
\subsection{Dummy Variable Trap}
\label{subsec:dummy-variable-trap}
%%%%%%%%%%%%%%%%%%%%%%%%%%%%%%%%%%%%%%%%%%%%%%%%%%%%%%%%%%%%

Suppose

\[
D_1+\cdots+D_k=1.
\]

If all \(k\) category indicators and an intercept are included, then one
column of the design matrix is an exact linear combination of the
others.

Thus

\[
X^TX
\]

is singular.

This problem is sometimes called the dummy-variable trap.

%%%%%%%%%%%%%%%%%%%%%%%%%%%%%%%%%%%%%%%%%%%%%%%%%%%%%%%%%%%%
\subsection{Interaction Terms}
\label{subsec:interaction-terms}
%%%%%%%%%%%%%%%%%%%%%%%%%%%%%%%%%%%%%%%%%%%%%%%%%%%%%%%%%%%%

Suppose the model is

\[
\boxed{
Y
=
\beta_0
+
\beta_1X
+
\beta_2Z
+
\beta_3XZ
+
\varepsilon.
}
\]

The effect of \(X\) depends on \(Z\):

\[
\frac{
\partial
\mathbb{E}[Y\mid X,Z]
}{
\partial X
}
=
\boxed{
\beta_1+\beta_3Z.
}
\]

Thus

\[
\beta_3
\]

measures interaction between \(X\) and \(Z\).

%%%%%%%%%%%%%%%%%%%%%%%%%%%%%%%%%%%%%%%%%%%%%%%%%%%%%%%%%%%%
\subsection{Worked Example: Interaction Interpretation}
\label{subsec:worked-interaction}
%%%%%%%%%%%%%%%%%%%%%%%%%%%%%%%%%%%%%%%%%%%%%%%%%%%%%%%%%%%%

\begin{workedexamplebox}{Different Slopes by Group}

Consider

\[
Y
=
10
+
2X
+
5D
-
X D
+
\varepsilon,
\]

where \(D=0\) or \(1\).

For \(D=0\),

\[
\mathbb{E}[Y\mid X,D=0]
=
10+2X.
\]

For \(D=1\),

\[
\mathbb{E}[Y\mid X,D=1]
=
15+X.
\]

Thus,

\begin{answerbox}
\[
\boxed{
\text{slope for }D=0:2,
\qquad
\text{slope for }D=1:1.
}
\]
\end{answerbox}

\end{workedexamplebox}

%%%%%%%%%%%%%%%%%%%%%%%%%%%%%%%%%%%%%%%%%%%%%%%%%%%%%%%%%%%%
\subsection{Polynomial Regression}
\label{subsec:polynomial-regression}
%%%%%%%%%%%%%%%%%%%%%%%%%%%%%%%%%%%%%%%%%%%%%%%%%%%%%%%%%%%%

A quadratic regression model is

\[
\boxed{
Y
=
\beta_0
+
\beta_1X
+
\beta_2X^2
+
\varepsilon.
}
\]

This model is nonlinear in \(X\) but linear in the coefficients

\[
\beta_0,\beta_1,\beta_2.
\]

Therefore ordinary linear-model methods still apply.

%%%%%%%%%%%%%%%%%%%%%%%%%%%%%%%%%%%%%%%%%%%%%%%%%%%%%%%%%%%%
\subsection{Higher-Order Polynomial Regression}
\label{subsec:higher-order-polynomial-regression}
%%%%%%%%%%%%%%%%%%%%%%%%%%%%%%%%%%%%%%%%%%%%%%%%%%%%%%%%%%%%

A degree-\(k\) polynomial regression has

\[
\boxed{
Y
=
\beta_0
+
\beta_1X
+
\cdots
+
\beta_kX^k
+
\varepsilon.
}
\]

High-degree polynomials can fit complex shapes but may oscillate
strongly and extrapolate poorly.

Flexible alternatives such as splines are often preferable.

%%%%%%%%%%%%%%%%%%%%%%%%%%%%%%%%%%%%%%%%%%%%%%%%%%%%%%%%%%%%
\subsection{Transformations}
\label{subsec:regression-transformations}
%%%%%%%%%%%%%%%%%%%%%%%%%%%%%%%%%%%%%%%%%%%%%%%%%%%%%%%%%%%%

Transformations may improve linearity or variance behavior.

Examples include

\[
\boxed{
\log Y,
}
\]

\[
\boxed{
\sqrt{Y},
}
\]

\[
\boxed{
\log X,
}
\]

or

\[
\boxed{
1/X.
}
\]

Interpretation must be adjusted to the transformed scale.

%%%%%%%%%%%%%%%%%%%%%%%%%%%%%%%%%%%%%%%%%%%%%%%%%%%%%%%%%%%%
\subsection{Log-Linear Regression}
\label{subsec:log-linear-regression}
%%%%%%%%%%%%%%%%%%%%%%%%%%%%%%%%%%%%%%%%%%%%%%%%%%%%%%%%%%%%

Consider

\[
\boxed{
\log Y
=
\beta_0+\beta_1X+\varepsilon.
}
\]

Then

\[
Y
\]

changes multiplicatively as \(X\) changes.

A one-unit increase in \(X\) multiplies the modeled geometric scale by

\[
\boxed{
e^{\beta_1}.
}
\]

For small \(\beta_1\),

\[
100\beta_1\%
\]

is approximately the percentage change in \(Y\), under the appropriate
interpretation.

%%%%%%%%%%%%%%%%%%%%%%%%%%%%%%%%%%%%%%%%%%%%%%%%%%%%%%%%%%%%
\subsection{Log-Log Regression}
\label{subsec:log-log-regression}
%%%%%%%%%%%%%%%%%%%%%%%%%%%%%%%%%%%%%%%%%%%%%%%%%%%%%%%%%%%%

Consider

\[
\boxed{
\log Y
=
\beta_0
+
\beta_1\log X
+
\varepsilon.
}
\]

Then

\[
\beta_1
\]

is an elasticity.

A \(1\%\) increase in \(X\) is associated with approximately a

\[
\boxed{
\beta_1\%
}
\]

change in the modeled conditional response scale.

%%%%%%%%%%%%%%%%%%%%%%%%%%%%%%%%%%%%%%%%%%%%%%%%%%%%%%%%%%%%
\subsection{Residual Diagnostics}
\label{subsec:residual-diagnostics}
%%%%%%%%%%%%%%%%%%%%%%%%%%%%%%%%%%%%%%%%%%%%%%%%%%%%%%%%%%%%

Residual diagnostics evaluate whether the fitted model is compatible with
its assumptions.

Useful plots include:

\[
\boxed{
\text{residuals versus fitted values},
}
\]

\[
\boxed{
\text{residuals versus predictors},
}
\]

\[
\boxed{
\text{normal Q--Q plot},
}
\]

and

\[
\boxed{
\text{scale-location plots}.
}
\]

Patterns in residuals suggest model deficiencies.

%%%%%%%%%%%%%%%%%%%%%%%%%%%%%%%%%%%%%%%%%%%%%%%%%%%%%%%%%%%%
\subsection{Residuals Versus Fitted Values}
\label{subsec:residuals-versus-fitted}
%%%%%%%%%%%%%%%%%%%%%%%%%%%%%%%%%%%%%%%%%%%%%%%%%%%%%%%%%%%%

A well-specified simple regression often produces residuals scattered
around zero without systematic structure.

A curved pattern may indicate:

\[
\boxed{
\text{nonlinearity}.
}
\]

A funnel shape may indicate:

\[
\boxed{
\text{heteroskedasticity}.
}
\]

Clusters may indicate missing group structure.

%%%%%%%%%%%%%%%%%%%%%%%%%%%%%%%%%%%%%%%%%%%%%%%%%%%%%%%%%%%%
\subsection{Heteroskedasticity}
\label{subsec:heteroskedasticity}
%%%%%%%%%%%%%%%%%%%%%%%%%%%%%%%%%%%%%%%%%%%%%%%%%%%%%%%%%%%%

Heteroskedasticity means

\[
\boxed{
\operatorname{Var}
(
\varepsilon_i\mid X_i
)
}
\]

is not constant.

OLS coefficient estimates can remain unbiased under zero conditional
mean, but the usual homoskedastic standard-error formulas can be
incorrect.

Thus the main problem is often inference rather than point estimation.

%%%%%%%%%%%%%%%%%%%%%%%%%%%%%%%%%%%%%%%%%%%%%%%%%%%%%%%%%%%%
\subsection{Heteroskedasticity-Robust Standard Errors}
\label{subsec:robust-standard-errors}
%%%%%%%%%%%%%%%%%%%%%%%%%%%%%%%%%%%%%%%%%%%%%%%%%%%%%%%%%%%%

Heteroskedasticity-consistent covariance estimators replace the
homoskedastic covariance formula

\[
\sigma^2(X^TX)^{-1}
\]

with a sandwich estimator.

A schematic form is

\[
\boxed{
(X^TX)^{-1}
X^T
\widehat{\Omega}
X
(X^TX)^{-1}.
}
\]

Different HC corrections modify the estimate of
\(\widehat{\Omega}\).

These methods leave OLS coefficients unchanged but adjust their estimated
uncertainty.

%%%%%%%%%%%%%%%%%%%%%%%%%%%%%%%%%%%%%%%%%%%%%%%%%%%%%%%%%%%%
\subsection{Weighted Least Squares}
\label{subsec:weighted-least-squares}
%%%%%%%%%%%%%%%%%%%%%%%%%%%%%%%%%%%%%%%%%%%%%%%%%%%%%%%%%%%%

If error variances differ in a known or modeled way, weighted least
squares can improve efficiency.

Minimize

\[
\boxed{
\sum_{i=1}^{n}
w_i
(y_i-\widehat{y}_i)^2.
}
\]

Observations with smaller error variance typically receive larger
weights.

In matrix form,

\[
\boxed{
\widehat{\boldsymbol{\beta}}_{WLS}
=
(X^TWX)^{-1}
X^TW\mathbf{y}.
}
\]

%%%%%%%%%%%%%%%%%%%%%%%%%%%%%%%%%%%%%%%%%%%%%%%%%%%%%%%%%%%%
\subsection{Generalized Least Squares}
\label{subsec:generalized-least-squares}
%%%%%%%%%%%%%%%%%%%%%%%%%%%%%%%%%%%%%%%%%%%%%%%%%%%%%%%%%%%%

Suppose

\[
\operatorname{Var}
(
\boldsymbol{\varepsilon}
)
=
\Sigma
\]

rather than \(\sigma^2I\).

If \(\Sigma\) is known and nonsingular, generalized least squares uses

\[
\boxed{
\widehat{\boldsymbol{\beta}}_{GLS}
=
(X^T\Sigma^{-1}X)^{-1}
X^T\Sigma^{-1}\mathbf{y}.
}
\]

This accounts for heteroskedasticity or correlation in the errors.

%%%%%%%%%%%%%%%%%%%%%%%%%%%%%%%%%%%%%%%%%%%%%%%%%%%%%%%%%%%%
\subsection{Normal Q--Q Plot}
\label{subsec:normal-qq-regression}
%%%%%%%%%%%%%%%%%%%%%%%%%%%%%%%%%%%%%%%%%%%%%%%%%%%%%%%%%%%%

A normal Q--Q plot compares ordered residuals with theoretical normal
quantiles.

Points approximately following a straight line suggest normality may be
reasonable.

Strong curvature or extreme tail deviations suggest nonnormality.

Normal Q--Q plots are more informative than simply inspecting a
histogram for small datasets.

%%%%%%%%%%%%%%%%%%%%%%%%%%%%%%%%%%%%%%%%%%%%%%%%%%%%%%%%%%%%
\subsection{Outliers}
\label{subsec:regression-outliers}
%%%%%%%%%%%%%%%%%%%%%%%%%%%%%%%%%%%%%%%%%%%%%%%%%%%%%%%%%%%%

A regression outlier is an observation with an unusual response relative
to the fitted model.

It may have a large residual.

However, an observation can have a large residual without having unusual
predictor values.

Thus outlyingness in the response and leverage in the predictors are
different concepts.

%%%%%%%%%%%%%%%%%%%%%%%%%%%%%%%%%%%%%%%%%%%%%%%%%%%%%%%%%%%%
\subsection{Leverage}
\label{subsec:regression-leverage}
%%%%%%%%%%%%%%%%%%%%%%%%%%%%%%%%%%%%%%%%%%%%%%%%%%%%%%%%%%%%

The diagonal elements

\[
h_{ii}
\]

of the hat matrix measure leverage.

Since

\[
\widehat{\mathbf{y}}
=
H\mathbf{y},
\]

the value \(h_{ii}\) reflects how strongly observation \(i\)'s fitted
value depends on its own observed response.

Large leverage generally occurs for unusual predictor combinations.

%%%%%%%%%%%%%%%%%%%%%%%%%%%%%%%%%%%%%%%%%%%%%%%%%%%%%%%%%%%%
\subsection{Average Leverage}
\label{subsec:average-leverage}
%%%%%%%%%%%%%%%%%%%%%%%%%%%%%%%%%%%%%%%%%%%%%%%%%%%%%%%%%%%%

Because

\[
\operatorname{tr}(H)=p+1,
\]

we have

\[
\sum_{i=1}^{n}
h_{ii}
=
p+1.
\]

Therefore the average leverage is

\[
\boxed{
\frac{
p+1
}{n}.
}
\]

Values much larger than this average may deserve investigation.

%%%%%%%%%%%%%%%%%%%%%%%%%%%%%%%%%%%%%%%%%%%%%%%%%%%%%%%%%%%%
\subsection{Studentized Residuals}
\label{subsec:studentized-residuals}
%%%%%%%%%%%%%%%%%%%%%%%%%%%%%%%%%%%%%%%%%%%%%%%%%%%%%%%%%%%%

Residual variances differ because

\[
\operatorname{Var}(e_i)
=
\sigma^2(1-h_{ii}).
\]

Therefore a standardized residual uses

\[
\boxed{
r_i
=
\frac{
e_i
}{
S\sqrt{1-h_{ii}}
}.
}
\]

Studentized residuals make residual magnitudes more comparable across
observations.

%%%%%%%%%%%%%%%%%%%%%%%%%%%%%%%%%%%%%%%%%%%%%%%%%%%%%%%%%%%%
\subsection{Influence}
\label{subsec:regression-influence}
%%%%%%%%%%%%%%%%%%%%%%%%%%%%%%%%%%%%%%%%%%%%%%%%%%%%%%%%%%%%

An observation is influential if removing it substantially changes the
fitted model.

Influence depends jointly on:

\[
\boxed{
\text{residual size}
}
\]

and

\[
\boxed{
\text{leverage}.
}
\]

A high-leverage point lying exactly on the regression trend may have
little influence.

A moderate residual combined with extreme leverage can be highly
influential.

%%%%%%%%%%%%%%%%%%%%%%%%%%%%%%%%%%%%%%%%%%%%%%%%%%%%%%%%%%%%
\subsection{Cook's Distance}
\label{subsec:cooks-distance}
%%%%%%%%%%%%%%%%%%%%%%%%%%%%%%%%%%%%%%%%%%%%%%%%%%%%%%%%%%%%

Cook's distance measures how much fitted values change when observation
\(i\) is removed.

A common form is

\[
\boxed{
D_i
=
\frac{
e_i^2
}{
(p+1)S^2
}
\frac{
h_{ii}
}{
(1-h_{ii})^2
}.
}
\]

Large values indicate observations with substantial potential influence.

Thresholds should be used as diagnostics rather than automatic deletion
rules.

%%%%%%%%%%%%%%%%%%%%%%%%%%%%%%%%%%%%%%%%%%%%%%%%%%%%%%%%%%%%
\subsection{Outliers Should Not Be Deleted Automatically}
\label{subsec:outlier-deletion-regression}
%%%%%%%%%%%%%%%%%%%%%%%%%%%%%%%%%%%%%%%%%%%%%%%%%%%%%%%%%%%%

An unusual observation may represent:

\[
\boxed{
\text{data error},
}
\]

\[
\boxed{
\text{rare valid behavior},
}
\]

\[
\boxed{
\text{model misspecification},
}
\]

or

\[
\boxed{
\text{a scientifically important subgroup}.
}
\]

Investigation should precede exclusion.

%%%%%%%%%%%%%%%%%%%%%%%%%%%%%%%%%%%%%%%%%%%%%%%%%%%%%%%%%%%%
\subsection{Multicollinearity}
\label{subsec:multicollinearity}
%%%%%%%%%%%%%%%%%%%%%%%%%%%%%%%%%%%%%%%%%%%%%%%%%%%%%%%%%%%%

Multicollinearity occurs when explanatory variables are highly linearly
related.

Exact multicollinearity makes

\[
X^TX
\]

singular.

Near multicollinearity makes it poorly conditioned.

Consequences can include:

\[
\boxed{
\text{large coefficient standard errors},
}
\]

\[
\boxed{
\text{unstable coefficient signs},
}
\]

and

\[
\boxed{
\text{sensitivity to small data changes}.
}
\]

%%%%%%%%%%%%%%%%%%%%%%%%%%%%%%%%%%%%%%%%%%%%%%%%%%%%%%%%%%%%
\subsection{Variance Inflation Factor}
\label{subsec:variance-inflation-factor}
%%%%%%%%%%%%%%%%%%%%%%%%%%%%%%%%%%%%%%%%%%%%%%%%%%%%%%%%%%%%

For predictor \(X_j\), regress it on the other predictors and let the
resulting coefficient of determination be

\[
R_j^2.
\]

Then the variance inflation factor is

\[
\boxed{
VIF_j
=
\frac1{
1-R_j^2
}.
}
\]

If \(R_j^2\) is near \(1\), then \(VIF_j\) is large.

This indicates that \(X_j\) is highly predictable from the remaining
predictors.

%%%%%%%%%%%%%%%%%%%%%%%%%%%%%%%%%%%%%%%%%%%%%%%%%%%%%%%%%%%%
\subsection{Multicollinearity Does Not Necessarily Harm Prediction}
\label{subsec:multicollinearity-prediction}
%%%%%%%%%%%%%%%%%%%%%%%%%%%%%%%%%%%%%%%%%%%%%%%%%%%%%%%%%%%%

Severe multicollinearity can make individual coefficient estimates
unstable.

However, predictions within the observed predictor region may still be
accurate.

Thus:

\[
\boxed{
\text{coefficient interpretation}
}
\]

and

\[
\boxed{
\text{prediction accuracy}
}
\]

are distinct goals.

%%%%%%%%%%%%%%%%%%%%%%%%%%%%%%%%%%%%%%%%%%%%%%%%%%%%%%%%%%%%
\subsection{Omitted Variable Bias}
\label{subsec:omitted-variable-bias}
%%%%%%%%%%%%%%%%%%%%%%%%%%%%%%%%%%%%%%%%%%%%%%%%%%%%%%%%%%%%

Suppose the true model contains

\[
Y
=
\beta_0+\beta_1X+\beta_2Z+\varepsilon,
\]

but \(Z\) is omitted.

If \(Z\) affects \(Y\) and is correlated with \(X\), the estimated
coefficient on \(X\) may absorb part of \(Z\)'s effect.

Thus the slope can be systematically distorted.

This is called omitted-variable bias.

%%%%%%%%%%%%%%%%%%%%%%%%%%%%%%%%%%%%%%%%%%%%%%%%%%%%%%%%%%%%
\subsection{Regression and Causation}
\label{subsec:regression-causation}
%%%%%%%%%%%%%%%%%%%%%%%%%%%%%%%%%%%%%%%%%%%%%%%%%%%%%%%%%%%%

A regression coefficient measures conditional association under the
specified model.

It is not automatically a causal effect.

Causal interpretation may require assumptions involving:

\[
\boxed{
\text{randomization},
}
\]

\[
\boxed{
\text{no unmeasured confounding},
}
\]

\[
\boxed{
\text{correct temporal ordering},
}
\]

and

\[
\boxed{
\text{appropriate treatment definition}.
}
\]

%%%%%%%%%%%%%%%%%%%%%%%%%%%%%%%%%%%%%%%%%%%%%%%%%%%%%%%%%%%%
\subsection{Simpson's Paradox}
\label{subsec:simpsons-paradox-regression}
%%%%%%%%%%%%%%%%%%%%%%%%%%%%%%%%%%%%%%%%%%%%%%%%%%%%%%%%%%%%

An association observed in aggregated data can reverse after conditioning
on another variable.

This phenomenon is known as Simpson's paradox.

Regression can help expose such relationships by adjusting for
additional variables, but adjustment must be scientifically justified.

Conditioning on inappropriate variables can itself create bias.

%%%%%%%%%%%%%%%%%%%%%%%%%%%%%%%%%%%%%%%%%%%%%%%%%%%%%%%%%%%%
\subsection{Confounders Versus Mediators}
\label{subsec:confounders-mediators}
%%%%%%%%%%%%%%%%%%%%%%%%%%%%%%%%%%%%%%%%%%%%%%%%%%%%%%%%%%%%

A confounder influences both an explanatory variable and the outcome.

A mediator lies on a causal pathway from exposure to outcome.

Adjusting for a confounder may reduce bias.

Adjusting for a mediator changes the causal quantity being estimated.

Thus variable selection for causal regression cannot be based only on
statistical significance.

%%%%%%%%%%%%%%%%%%%%%%%%%%%%%%%%%%%%%%%%%%%%%%%%%%%%%%%%%%%%
\subsection{Model Selection}
\label{subsec:regression-model-selection}
%%%%%%%%%%%%%%%%%%%%%%%%%%%%%%%%%%%%%%%%%%%%%%%%%%%%%%%%%%%%

Model selection determines which explanatory variables or functional
forms to include.

Possible criteria include:

\[
\boxed{
\text{scientific theory},
}
\]

\[
\boxed{
\text{prediction performance},
}
\]

\[
\boxed{
\text{AIC},
}
\]

\[
\boxed{
\text{BIC},
}
\]

\[
\boxed{
\text{cross-validation},
}
\]

and

\[
\boxed{
\text{regularization}.
}
\]

The best procedure depends on the goal.

%%%%%%%%%%%%%%%%%%%%%%%%%%%%%%%%%%%%%%%%%%%%%%%%%%%%%%%%%%%%
\subsection{Adjusted \(R^2\) Is Not a Universal Model Selector}
\label{subsec:adjusted-r2-not-universal}
%%%%%%%%%%%%%%%%%%%%%%%%%%%%%%%%%%%%%%%%%%%%%%%%%%%%%%%%%%%%

Adjusted \(R^2\) penalizes adding predictors, but it does not directly
measure:

\[
\boxed{
\text{out-of-sample prediction error},
}
\]

\[
\boxed{
\text{causal validity},
}
\]

or

\[
\boxed{
\text{model calibration}.
}
\]

Therefore model selection should not rely on a single statistic.

%%%%%%%%%%%%%%%%%%%%%%%%%%%%%%%%%%%%%%%%%%%%%%%%%%%%%%%%%%%%
\subsection{Akaike Information Criterion Preview}
\label{subsec:aic-preview}
%%%%%%%%%%%%%%%%%%%%%%%%%%%%%%%%%%%%%%%%%%%%%%%%%%%%%%%%%%%%

For a likelihood-based model with \(k\) estimated parameters,

\[
\boxed{
AIC
=
-2\ell(\widehat{\theta})
+
2k.
}
\]

Smaller AIC values indicate a better tradeoff between fit and model
complexity within the candidate model set.

AIC is motivated by predictive information loss rather than by
hypothesis-testing significance.

%%%%%%%%%%%%%%%%%%%%%%%%%%%%%%%%%%%%%%%%%%%%%%%%%%%%%%%%%%%%
\subsection{Bayesian Information Criterion Preview}
\label{subsec:bic-preview}
%%%%%%%%%%%%%%%%%%%%%%%%%%%%%%%%%%%%%%%%%%%%%%%%%%%%%%%%%%%%

The Bayesian information criterion is

\[
\boxed{
BIC
=
-2\ell(\widehat{\theta})
+
k\ln n.
}
\]

BIC penalizes model complexity more strongly than AIC as sample size
increases.

Its interpretation differs from that of AIC despite their similar form.

%%%%%%%%%%%%%%%%%%%%%%%%%%%%%%%%%%%%%%%%%%%%%%%%%%%%%%%%%%%%
\subsection{Cross-Validation Preview}
\label{subsec:cross-validation-regression}
%%%%%%%%%%%%%%%%%%%%%%%%%%%%%%%%%%%%%%%%%%%%%%%%%%%%%%%%%%%%

Cross-validation estimates out-of-sample predictive performance.

The dataset is repeatedly divided into training and validation portions.

A model is fit on training data and evaluated on held-out data.

This directly addresses prediction rather than merely in-sample fit.

%%%%%%%%%%%%%%%%%%%%%%%%%%%%%%%%%%%%%%%%%%%%%%%%%%%%%%%%%%%%
\subsection{Overfitting}
\label{subsec:regression-overfitting}
%%%%%%%%%%%%%%%%%%%%%%%%%%%%%%%%%%%%%%%%%%%%%%%%%%%%%%%%%%%%

Overfitting occurs when a model adapts too closely to random noise in the
training data.

An overfit model may have:

\[
\boxed{
\text{excellent in-sample fit}
}
\]

but

\[
\boxed{
\text{poor out-of-sample prediction}.
}
\]

Adding predictors always reduces or leaves unchanged \(SSE\), so
training fit alone cannot detect overfitting.

%%%%%%%%%%%%%%%%%%%%%%%%%%%%%%%%%%%%%%%%%%%%%%%%%%%%%%%%%%%%
\subsection{Underfitting}
\label{subsec:regression-underfitting}
%%%%%%%%%%%%%%%%%%%%%%%%%%%%%%%%%%%%%%%%%%%%%%%%%%%%%%%%%%%%

Underfitting occurs when a model is too simple to capture important
structure.

Examples include fitting a straight line to a strongly curved
relationship.

Thus model complexity involves a tradeoff:

\[
\boxed{
\text{too simple}
\longleftrightarrow
\text{too flexible}.
}
\]

%%%%%%%%%%%%%%%%%%%%%%%%%%%%%%%%%%%%%%%%%%%%%%%%%%%%%%%%%%%%
\subsection{Ridge Regression}
\label{subsec:ridge-regression}
%%%%%%%%%%%%%%%%%%%%%%%%%%%%%%%%%%%%%%%%%%%%%%%%%%%%%%%%%%%%

Ridge regression minimizes

\[
\boxed{
\|\mathbf{y}-X\boldsymbol{\beta}\|_2^2
+
\lambda
\sum_{j=1}^{p}
\beta_j^2.
}
\]

Usually the intercept is not penalized.

The penalty shrinks coefficients toward zero.

For appropriately standardized predictors,

\[
\boxed{
\widehat{\boldsymbol{\beta}}_{\text{ridge}}
=
(X^TX+\lambda I)^{-1}
X^T\mathbf{y}
}
\]

for the penalized coefficient block under the relevant convention.

%%%%%%%%%%%%%%%%%%%%%%%%%%%%%%%%%%%%%%%%%%%%%%%%%%%%%%%%%%%%
\subsection{Why Ridge Regression Helps}
\label{subsec:why-ridge}
%%%%%%%%%%%%%%%%%%%%%%%%%%%%%%%%%%%%%%%%%%%%%%%%%%%%%%%%%%%%

Ridge regression introduces bias but can reduce variance.

It is particularly useful when:

\[
\boxed{
\text{predictors are highly correlated},
}
\]

or

\[
\boxed{
p
\text{ is large relative to }n.
}
\]

This is an example of the bias--variance tradeoff.

%%%%%%%%%%%%%%%%%%%%%%%%%%%%%%%%%%%%%%%%%%%%%%%%%%%%%%%%%%%%
\subsection{Lasso Regression}
\label{subsec:lasso-regression}
%%%%%%%%%%%%%%%%%%%%%%%%%%%%%%%%%%%%%%%%%%%%%%%%%%%%%%%%%%%%

The lasso minimizes

\[
\boxed{
\|\mathbf{y}-X\boldsymbol{\beta}\|_2^2
+
\lambda
\sum_{j=1}^{p}
|\beta_j|.
}
\]

The \(L^1\) penalty can force some estimated coefficients exactly to
zero.

Thus lasso can perform both:

\[
\boxed{
\text{shrinkage}
}
\]

and

\[
\boxed{
\text{variable selection}.
}
\]

%%%%%%%%%%%%%%%%%%%%%%%%%%%%%%%%%%%%%%%%%%%%%%%%%%%%%%%%%%%%
\subsection{Standardization Before Regularization}
\label{subsec:standardization-regularization}
%%%%%%%%%%%%%%%%%%%%%%%%%%%%%%%%%%%%%%%%%%%%%%%%%%%%%%%%%%%%

Because ridge and lasso penalties depend on coefficient magnitude,
predictor scale matters.

Therefore predictors are often standardized before fitting a regularized
model.

Otherwise variables measured in different units may be penalized
unequally for purely numerical reasons.

%%%%%%%%%%%%%%%%%%%%%%%%%%%%%%%%%%%%%%%%%%%%%%%%%%%%%%%%%%%%
\subsection{Generalized Linear Models Preview}
\label{subsec:glm-preview}
%%%%%%%%%%%%%%%%%%%%%%%%%%%%%%%%%%%%%%%%%%%%%%%%%%%%%%%%%%%%

Linear regression assumes a continuous response with conditional mean
modeled directly by a linear predictor.

Generalized linear models extend regression to other response
distributions.

They use:

\[
\boxed{
\text{random component}
+
\text{linear predictor}
+
\text{link function}.
}
\]

Examples include logistic and Poisson regression.

%%%%%%%%%%%%%%%%%%%%%%%%%%%%%%%%%%%%%%%%%%%%%%%%%%%%%%%%%%%%
\subsection{Logistic Regression Preview}
\label{subsec:logistic-regression-preview}
%%%%%%%%%%%%%%%%%%%%%%%%%%%%%%%%%%%%%%%%%%%%%%%%%%%%%%%%%%%%

For a binary response,

\[
Y\in\{0,1\},
\]

let

\[
p(x)
=
P(Y=1\mid X=x).
\]

Logistic regression models the log-odds:

\[
\boxed{
\log
\left(
\frac{
p(x)
}{
1-p(x)
}
\right)
=
\beta_0+\beta_1x.
}
\]

Solving for \(p(x)\),

\[
\boxed{
p(x)
=
\frac{
e^{\beta_0+\beta_1x}
}{
1+
e^{\beta_0+\beta_1x}
}.
}
\]

%%%%%%%%%%%%%%%%%%%%%%%%%%%%%%%%%%%%%%%%%%%%%%%%%%%%%%%%%%%%
\subsection{Odds-Ratio Interpretation Preview}
\label{subsec:odds-ratio-logistic-preview}
%%%%%%%%%%%%%%%%%%%%%%%%%%%%%%%%%%%%%%%%%%%%%%%%%%%%%%%%%%%%

In logistic regression,

\[
\beta_1
\]

is the change in log-odds for a one-unit increase in \(X\).

Exponentiating,

\[
\boxed{
e^{\beta_1}
}
\]

is the multiplicative change in the odds for a one-unit increase in
\(X\), holding other predictors fixed.

%%%%%%%%%%%%%%%%%%%%%%%%%%%%%%%%%%%%%%%%%%%%%%%%%%%%%%%%%%%%
\subsection{Poisson Regression Preview}
\label{subsec:poisson-regression-preview}
%%%%%%%%%%%%%%%%%%%%%%%%%%%%%%%%%%%%%%%%%%%%%%%%%%%%%%%%%%%%

For count responses, one may model

\[
Y_i
\sim
\operatorname{Poisson}(\mu_i)
\]

with

\[
\boxed{
\log\mu_i
=
\beta_0
+
\beta_1x_{i1}
+
\cdots
+
\beta_px_{ip}.
}
\]

Thus,

\[
\boxed{
\mu_i
=
\exp
(
\beta_0+\cdots+\beta_px_{ip}
).
}
\]

This guarantees a positive predicted mean.

%%%%%%%%%%%%%%%%%%%%%%%%%%%%%%%%%%%%%%%%%%%%%%%%%%%%%%%%%%%%
\subsection{Regression Modeling Workflow}
\label{subsec:regression-modeling-workflow}
%%%%%%%%%%%%%%%%%%%%%%%%%%%%%%%%%%%%%%%%%%%%%%%%%%%%%%%%%%%%

A practical regression workflow is:

\begin{enumerate}
    \item identify the response variable;
    \item identify candidate explanatory variables;
    \item inspect data and scatterplots;
    \item define the scientific or predictive goal;
    \item specify a model;
    \item fit the model;
    \item interpret coefficients;
    \item examine standard errors and confidence intervals;
    \item test scientifically meaningful hypotheses;
    \item inspect residual diagnostics;
    \item assess leverage and influence;
    \item investigate collinearity;
    \item evaluate predictive performance;
    \item check sensitivity to modeling choices;
    \item avoid causal interpretations not supported by the study design.
\end{enumerate}

%%%%%%%%%%%%%%%%%%%%%%%%%%%%%%%%%%%%%%%%%%%%%%%%%%%%%%%%%%%%
\subsection{Common Errors}
\label{subsec:regression-common-errors}
%%%%%%%%%%%%%%%%%%%%%%%%%%%%%%%%%%%%%%%%%%%%%%%%%%%%%%%%%%%%

\subsubsection{Interpreting Association as Causation}

A regression slope is not automatically a causal effect.

Causal interpretation depends on study design and assumptions.

%%%%%%%%%%%%%%%%%%%%%%%%%%%%%%%%%%%%%%%%%%%%%%%%%%%%%%%%%%%%
\subsubsection{Ignoring the Meaning of ``Holding Other Variables Fixed''}

A multiple-regression coefficient is a partial association conditional
on the other included variables.

It is not generally the same as the simple bivariate relationship.

%%%%%%%%%%%%%%%%%%%%%%%%%%%%%%%%%%%%%%%%%%%%%%%%%%%%%%%%%%%%
\subsubsection{Interpreting a Nonsignificant Coefficient as No Relationship}

A large standard error may reflect limited data or multicollinearity.

Failure to reject

\[
\beta_j=0
\]

does not prove the coefficient is exactly zero.

%%%%%%%%%%%%%%%%%%%%%%%%%%%%%%%%%%%%%%%%%%%%%%%%%%%%%%%%%%%%
\subsubsection{Interpreting a Significant Coefficient as Important}

A very small coefficient can be statistically significant in a large
sample.

Magnitude and uncertainty should be reported together.

%%%%%%%%%%%%%%%%%%%%%%%%%%%%%%%%%%%%%%%%%%%%%%%%%%%%%%%%%%%%
\subsubsection{Using \(R^2\) as the Only Measure of Model Quality}

High \(R^2\) does not guarantee:

\[
\boxed{
\text{correct specification},
}
\]

\[
\boxed{
\text{good causal interpretation},
}
\]

or

\[
\boxed{
\text{good future prediction}.
}
\]

%%%%%%%%%%%%%%%%%%%%%%%%%%%%%%%%%%%%%%%%%%%%%%%%%%%%%%%%%%%%
\subsubsection{Comparing \(R^2\) Across Unrelated Responses Without Context}

The natural amount of variation differs across scientific settings.

An \(R^2\) that is useful in one field may be inadequate or exceptional
in another.

%%%%%%%%%%%%%%%%%%%%%%%%%%%%%%%%%%%%%%%%%%%%%%%%%%%%%%%%%%%%
\subsubsection{Ignoring Residual Patterns}

A good numerical fit can still hide nonlinearity, heteroskedasticity, or
subgroup structure.

Residual diagnostics should accompany regression fitting.

%%%%%%%%%%%%%%%%%%%%%%%%%%%%%%%%%%%%%%%%%%%%%%%%%%%%%%%%%%%%
\subsubsection{Deleting High-Leverage Observations Automatically}

High leverage alone does not imply an observation is invalid.

It indicates unusual predictor values.

%%%%%%%%%%%%%%%%%%%%%%%%%%%%%%%%%%%%%%%%%%%%%%%%%%%%%%%%%%%%
\subsubsection{Deleting Large Residuals Automatically}

Large residuals can reveal important scientific structure or model
failure.

Investigate before excluding observations.

%%%%%%%%%%%%%%%%%%%%%%%%%%%%%%%%%%%%%%%%%%%%%%%%%%%%%%%%%%%%
\subsubsection{Ignoring Multicollinearity}

Highly related predictors can make coefficient estimates unstable even
when overall prediction remains strong.

%%%%%%%%%%%%%%%%%%%%%%%%%%%%%%%%%%%%%%%%%%%%%%%%%%%%%%%%%%%%
\subsubsection{Using a Linear Model Outside the Data Range}

Extrapolation assumes the modeled relationship continues beyond observed
predictor values.

This assumption may be unjustified.

%%%%%%%%%%%%%%%%%%%%%%%%%%%%%%%%%%%%%%%%%%%%%%%%%%%%%%%%%%%%
\subsubsection{Confusing Prediction Intervals and Mean-Response Intervals}

Prediction intervals target individual future observations and are
wider.

Mean-response intervals target the conditional mean.

%%%%%%%%%%%%%%%%%%%%%%%%%%%%%%%%%%%%%%%%%%%%%%%%%%%%%%%%%%%%
\subsubsection{Using Classical Standard Errors Under Heteroskedasticity}

If error variance is not constant, homoskedastic OLS standard errors can
be incorrect.

Robust or model-based alternatives may be needed.

%%%%%%%%%%%%%%%%%%%%%%%%%%%%%%%%%%%%%%%%%%%%%%%%%%%%%%%%%%%%
\subsubsection{Ignoring Dependence}

Regression errors may be correlated in time, space, clusters, or repeated
measurements.

Independence-based inference can then underestimate uncertainty.

%%%%%%%%%%%%%%%%%%%%%%%%%%%%%%%%%%%%%%%%%%%%%%%%%%%%%%%%%%%%
\subsubsection{Adding Variables Solely Because They Are Significant}

Variable selection should reflect the analysis goal and scientific
structure.

Pure significance-based stepwise procedures can produce unstable
estimates and invalid post-selection inference.

%%%%%%%%%%%%%%%%%%%%%%%%%%%%%%%%%%%%%%%%%%%%%%%%%%%%%%%%%%%%
\subsection{Applications}
\label{subsec:regression-applications}
%%%%%%%%%%%%%%%%%%%%%%%%%%%%%%%%%%%%%%%%%%%%%%%%%%%%%%%%%%%%

\begin{applicationbox}

\textbf{Scientific Research.}
Regression models quantify relationships among measured variables and
estimate adjusted associations.

\medskip

\textbf{Engineering.}
Regression predicts system outputs from operating conditions and
calibrates empirical models.

\medskip

\textbf{Environmental Science.}
Environmental responses may be modeled using temperature, time, depth,
location, chemical concentration, or other predictors.

\medskip

\textbf{Education Research.}
Regression models student outcomes using preparation, demographics,
program participation, and other measured variables.

\medskip

\textbf{Economics.}
Regression estimates associations among prices, income, employment,
policy variables, and economic outcomes.

\medskip

\textbf{Medicine and Epidemiology.}
Regression adjusts for multiple risk factors and models treatment and
exposure relationships.

\medskip

\textbf{Business.}
Regression predicts demand, costs, sales, conversion rates, and customer
behavior.

\medskip

\textbf{Machine Learning.}
Linear models, regularization, logistic regression, and generalized
linear models form the foundation of many predictive algorithms.

\end{applicationbox}

%%%%%%%%%%%%%%%%%%%%%%%%%%%%%%%%%%%%%%%%%%%%%%%%%%%%%%%%%%%%
\subsection{Exercises}
\label{subsec:regression-exercises}
%%%%%%%%%%%%%%%%%%%%%%%%%%%%%%%%%%%%%%%%%%%%%%%%%%%%%%%%%%%%

\begin{exercise}[1]
Define simple linear regression.
\end{exercise}

\begin{exercise}[1]
Define the response variable and explanatory variable.
\end{exercise}

\begin{exercise}[1]
Interpret \(\beta_0\) and \(\beta_1\).
\end{exercise}

\begin{exercise}[1]
Define a fitted value.
\end{exercise}

\begin{exercise}[1]
Define a residual.
\end{exercise}

\begin{exercise}[1]
State the ordinary least-squares criterion.
\end{exercise}

\begin{exercise}[1]
Define \(SST\), \(SSR\), and \(SSE\).
\end{exercise}

\begin{exercise}[1]
Define \(R^2\).
\end{exercise}

\begin{exercise}[1]
State the classical simple linear regression assumptions.
\end{exercise}

\begin{exercise}[1]
Define homoskedasticity.
\end{exercise}

\begin{exercise}[1]
Define heteroskedasticity.
\end{exercise}

\begin{exercise}[1]
Define leverage.
\end{exercise}

\begin{exercise}[1]
Define multicollinearity.
\end{exercise}

\begin{exercise}[1]
Define a variance inflation factor.
\end{exercise}

\begin{exercise}[2]
Suppose

\[
\widehat{y}
=
5+3x.
\]

Interpret the slope.
\end{exercise}

\begin{exercise}[2]
For the same model, predict \(Y\) at

\[
x=4.
\]
\end{exercise}

\begin{exercise}[2]
If the observed response at \(x=4\) is \(20\), compute the residual.
\end{exercise}

\begin{exercise}[2]
For data with

\[
S_{xy}=24
\]

and

\[
S_{xx}=8,
\]

compute the OLS slope.
\end{exercise}

\begin{exercise}[2]
Suppose

\[
\overline{x}=5,
\qquad
\overline{y}=20,
\qquad
b_1=2.
\]

Compute the intercept.
\end{exercise}

\begin{exercise}[2]
Suppose

\[
SST=100,
\qquad
SSE=25.
\]

Compute \(SSR\) and \(R^2\).
\end{exercise}

\begin{exercise}[2]
Suppose

\[
r=-0.8.
\]

Find \(R^2\) for simple linear regression with an intercept.
\end{exercise}

\begin{exercise}[2]
Suppose

\[
r=0.5,
\qquad
s_y=12,
\qquad
s_x=3.
\]

Compute the simple-regression slope.
\end{exercise}

\begin{exercise}[2]
Suppose

\[
SSE=90,
\qquad
n=12.
\]

Compute the residual standard error for simple regression.
\end{exercise}

\begin{exercise}[2]
Suppose

\[
b_1=1.5,
\qquad
SE(b_1)=0.5.
\]

Compute the \(t\)-statistic for testing

\[
H_0:\beta_1=0.
\]
\end{exercise}

\begin{exercise}[2]
Suppose a multiple regression has \(n=100\) observations and \(p=4\)
predictors. Find the residual degrees of freedom.
\end{exercise}

\begin{exercise}[2]
Suppose a categorical predictor has \(5\) categories. How many indicator
variables are needed with an intercept?
\end{exercise}

\begin{exercise}[2]
Suppose \(R_j^2=0.80\) when \(X_j\) is regressed on all remaining
predictors. Compute \(VIF_j\).
\end{exercise}

\begin{exercise}[3]
Derive the normal equations for simple linear regression.
\end{exercise}

\begin{exercise}[3]
Derive

\[
b_1
=
\frac{S_{xy}}{S_{xx}}.
\]
\end{exercise}

\begin{exercise}[3]
Derive

\[
b_0
=
\overline{y}
-
b_1\overline{x}.
\]
\end{exercise}

\begin{exercise}[3]
Prove that the fitted regression line passes through

\[
(\overline{x},\overline{y}).
\]
\end{exercise}

\begin{exercise}[3]
Prove that

\[
\sum e_i=0
\]

when an intercept is included.
\end{exercise}

\begin{exercise}[3]
Prove that

\[
\sum x_ie_i=0.
\]
\end{exercise}

\begin{exercise}[3]
Prove the decomposition

\[
SST
=
SSR
+
SSE.
\]
\end{exercise}

\begin{exercise}[3]
Show that

\[
R^2=r^2
\]

in simple linear regression with an intercept.
\end{exercise}

\begin{exercise}[3]
Show that

\[
b_1
=
r
\frac{s_y}{s_x}.
\]
\end{exercise}

\begin{exercise}[3]
Derive

\[
\operatorname{Var}(b_1)
=
\frac{
\sigma^2
}{
S_{xx}
}.
\]
\end{exercise}

\begin{exercise}[3]
Derive the variance of the intercept estimator.
\end{exercise}

\begin{exercise}[3]
Derive the covariance between the simple-regression slope and intercept.
\end{exercise}

\begin{exercise}[3]
Derive the standard error for the estimated mean response at \(x_0\).
\end{exercise}

\begin{exercise}[3]
Derive the standard error for predicting a future observation at \(x_0\).
\end{exercise}

\begin{exercise}[3]
Explain why prediction intervals are wider than mean-response confidence
intervals.
\end{exercise}

\begin{exercise}[3]
Derive the matrix normal equations

\[
X^TX\widehat{\boldsymbol{\beta}}
=
X^T\mathbf{y}.
\]
\end{exercise}

\begin{exercise}[3]
Derive

\[
\widehat{\boldsymbol{\beta}}
=
(X^TX)^{-1}X^T\mathbf{y}.
\]
\end{exercise}

\begin{exercise}[3]
Show that the hat matrix is symmetric.
\end{exercise}

\begin{exercise}[3]
Show that the hat matrix is idempotent.
\end{exercise}

\begin{exercise}[3]
Show that the residual vector is orthogonal to the column space of \(X\).
\end{exercise}

\begin{exercise}[3]
Show that

\[
\operatorname{Var}
(
\widehat{\boldsymbol{\beta}}\mid X
)
=
\sigma^2(X^TX)^{-1}.
\]
\end{exercise}

\begin{exercise}[3]
Explain why exact multicollinearity prevents unique OLS coefficient
estimation.
\end{exercise}

\begin{exercise}[3]
Interpret the interaction coefficient in

\[
Y
=
\beta_0
+
\beta_1X
+
\beta_2Z
+
\beta_3XZ
+
\varepsilon.
\]
\end{exercise}

\begin{exercise}[3]
Explain why polynomial regression remains a linear model in the
parameters.
\end{exercise}

\begin{exercise}[3]
Explain why heteroskedasticity changes standard errors without
necessarily biasing OLS coefficients.
\end{exercise}

\begin{exercise}[3]
Explain why an observation can have high leverage but a small residual.
\end{exercise}

\begin{exercise}[3]
Explain why high \(R^2\) does not establish causality.
\end{exercise}

\begin{exercise}[3]
Construct an omitted-variable example in which an observed regression
slope is misleading.
\end{exercise}

\begin{exercise}[3]
Explain how multicollinearity can inflate coefficient standard errors.
\end{exercise}

\begin{exercise}[3]
Explain why ridge regression can improve prediction under
multicollinearity.
\end{exercise}

\begin{exercise}[3]
Compare ridge and lasso penalties.
\end{exercise}

\begin{exercise}[4]
Prove the Gauss--Markov theorem.
\end{exercise}

\begin{exercise}[4]
Study the Frisch--Waugh--Lovell theorem and partial regression.
\end{exercise}

\begin{exercise}[4]
Derive the exact normal-theory distribution of the OLS coefficient
vector.
\end{exercise}

\begin{exercise}[4]
Derive the overall regression \(F\)-test from quadratic forms.
\end{exercise}

\begin{exercise}[4]
Study general linear hypotheses of the form

\[
H_0:
R\boldsymbol{\beta}
=
\mathbf{r}.
\]
\end{exercise}

\begin{exercise}[4]
Derive the corresponding partial \(F\)-test.
\end{exercise}

\begin{exercise}[4]
Study heteroskedasticity-consistent covariance estimators HC0--HC3.
\end{exercise}

\begin{exercise}[4]
Investigate the Breusch--Pagan test.
\end{exercise}

\begin{exercise}[4]
Investigate the White test for heteroskedasticity.
\end{exercise}

\begin{exercise}[4]
Study generalized least squares under correlated errors.
\end{exercise}

\begin{exercise}[4]
Investigate autocorrelation and regression with time-series errors.
\end{exercise}

\begin{exercise}[4]
Study Cook's distance, DFBETAS, and DFFITS.
\end{exercise}

\begin{exercise}[4]
Investigate regression splines and smoothing splines.
\end{exercise}

\begin{exercise}[4]
Study ridge regression in terms of singular-value decomposition.
\end{exercise}

\begin{exercise}[4]
Derive the lasso optimization problem and study its geometric
interpretation.
\end{exercise}

\begin{exercise}[4]
Investigate cross-validation for regression model selection.
\end{exercise}

\begin{exercise}[4]
Study post-selection inference.
\end{exercise}

\begin{exercise}[4]
Investigate logistic regression through maximum likelihood.
\end{exercise}

\begin{exercise}[4]
Study Poisson regression and generalized linear models.
\end{exercise}

%%%%%%%%%%%%%%%%%%%%%%%%%%%%%%%%%%%%%%%%%%%%%%%%%%%%%%%%%%%%
\subsection{Conceptual Summary}
\label{subsec:regression-summary}
%%%%%%%%%%%%%%%%%%%%%%%%%%%%%%%%%%%%%%%%%%%%%%%%%%%%%%%%%%%%

Simple linear regression models

\[
\boxed{
Y_i
=
\beta_0
+
\beta_1X_i
+
\varepsilon_i.
}
\]

The conditional mean is

\[
\boxed{
\mathbb{E}[Y\mid X]
=
\beta_0+\beta_1X.
}
\]

Ordinary least squares minimizes

\[
\boxed{
\sum e_i^2.
}
\]

The simple-regression slope is

\[
\boxed{
b_1
=
\frac{
\sum
(x_i-\overline{x})
(y_i-\overline{y})
}{
\sum
(x_i-\overline{x})^2
}.
}
\]

The intercept is

\[
\boxed{
b_0
=
\overline{y}
-
b_1\overline{x}.
}
\]

Variation decomposes as

\[
\boxed{
SST
=
SSR+SSE.
}
\]

The coefficient of determination is

\[
\boxed{
R^2
=
1-\frac{SSE}{SST}.
}
\]

For simple regression,

\[
\boxed{
R^2=r^2.
}
\]

In matrix form,

\[
\boxed{
\mathbf{y}
=
X\boldsymbol{\beta}
+
\boldsymbol{\varepsilon},
}
\]

and

\[
\boxed{
\widehat{\boldsymbol{\beta}}
=
(X^TX)^{-1}
X^T\mathbf{y}.
}
\]

Under homoskedastic independent errors,

\[
\boxed{
\operatorname{Var}
(
\widehat{\boldsymbol{\beta}}
\mid X
)
=
\sigma^2
(X^TX)^{-1}.
}
\]

Regression therefore combines

\[
\boxed{
\text{linear algebra}
+
\text{optimization}
+
\text{probability}
+
\text{estimation}
+
\text{hypothesis testing}.
}
\]

%%%%%%%%%%%%%%%%%%%%%%%%%%%%%%%%%%%%%%%%%%%%%%%%%%%%%%%%%%%%
\subsection{Section Connections}
\label{subsec:regression-connections}
%%%%%%%%%%%%%%%%%%%%%%%%%%%%%%%%%%%%%%%%%%%%%%%%%%%%%%%%%%%%

\chapterconnection{
Regression extends the estimation and testing framework developed in
Sections~\ref{sec:statistical-estimation},
\ref{sec:confidence-intervals}, and
\ref{sec:hypothesis-testing} to models in which a response depends on one
or more explanatory variables. The least-squares decomposition of
variation leads directly to the Analysis of Variance framework developed
in the next section. Multiple regression also provides a foundation for
experimental design, categorical-data modeling, generalized linear
models, time series, statistical learning, and high-dimensional
statistics. Linear Algebra supplies the projection interpretation of OLS,
while probability theory provides the sampling distributions needed for
coefficient inference.
}

%%%%%%%%%%%%%%%%%%%%%%%%%%%%%%%%%%%%%%%%%%%%%%%%%%%%%%%%%%%%
% SECTION SOURCE TRACKING
%%%%%%%%%%%%%%%%%%%%%%%%%%%%%%%%%%%%%%%%%%%%%%%%%%%%%%%%%%%%

%------------------------------------------------------------
% 8.7 Regression
%------------------------------------------------------------
%
% Source basis:
%   - Standard linear regression theory.
%   - Standard least-squares and linear-model theory.
%   - Standard regression diagnostics and model-selection concepts.
%
% Used for:
%   - simple linear regression;
%   - conditional mean models;
%   - slope and intercept interpretation;
%   - fitted values and residuals;
%   - ordinary least squares;
%   - normal equations;
%   - slope and intercept estimators;
%   - residual properties;
%   - SST, SSR, and SSE;
%   - R-squared;
%   - correlation-regression relationships;
%   - classical linear-model assumptions;
%   - unbiasedness of OLS;
%   - coefficient variances;
%   - residual variance;
%   - coefficient standard errors;
%   - t-tests and confidence intervals;
%   - confidence intervals for mean response;
%   - prediction intervals;
%   - multiple linear regression;
%   - matrix formulation;
%   - design matrices;
%   - OLS matrix estimator;
%   - hat matrix;
%   - projection interpretation;
%   - residual-maker matrix;
%   - covariance matrix of OLS;
%   - overall F-tests;
%   - regression ANOVA tables;
%   - adjusted R-squared;
%   - indicator variables;
%   - categorical predictors;
%   - interaction terms;
%   - polynomial regression;
%   - transformations;
%   - residual diagnostics;
%   - heteroskedasticity;
%   - robust standard errors;
%   - weighted and generalized least squares;
%   - leverage;
%   - studentized residuals;
%   - influence and Cook's distance;
%   - multicollinearity;
%   - variance inflation factors;
%   - omitted-variable bias;
%   - causal interpretation cautions;
%   - model selection;
%   - AIC and BIC previews;
%   - cross-validation;
%   - overfitting and underfitting;
%   - ridge and lasso regression;
%   - generalized linear model previews;
%   - logistic regression;
%   - Poisson regression;
%   - applications and exercises.
%
% Notes:
%   - Analysis of Variance is developed next in Section 8.8.
%   - Experimental Design follows in Section 8.9.
%   - Categorical Data Analysis later develops logistic and log-linear
%     models more fully.
%   - Statistical Learning develops regularization and predictive
%     modeling more deeply.
%
%%%%%%%%%%%%%%%%%%%%%%%%%%%%%%%%%%%%%%%%%%%%%%%%%%%%%%%%%%%%